% ============================================================
%  Projective Dynamic Logo — Popular science for general audiences
%  Cédric Laubscher — 2026
%  Compile with XeLaTeX or LuaLaTeX (recommended)
%  or pdfLaTeX (replace fontspec with inputenc/fontenc)
% ============================================================

\documentclass[12pt, a4paper, openright]{book}
\usepackage{hyphenat}
\usepackage{hyphenat}
% --- Encoding and language ---
\ifdefined\XeTeXversion
\usepackage{ifxetex,ifluatex}
\ifxetex
\usepackage{fontspec}
  \setmainfont{Latin Modern}           % replace with EB Garamond if unavailable
  \setsansfont{Source Sans Pro}
  \setmonofont{JuliaMono}[Scale=0.85]
\else
  \usepackage[T1]{fontenc}
  \usepackage[utf8]{inputenc}
  \usepackage{lmodern}
\fi

  \usepackage[utf8]{inputenc}
  \usepackage[T1]{fontenc}

\let\showhyphens\relax
\usepackage[english]{babel}
\usepackage{csquotes}

% --- Page layout ---
\usepackage{xcolor}
\usepackage[
  top=2.8cm, bottom=3cm,
  left=3cm, right=3cm,
  headheight=15pt
]{geometry}
\usepackage{emptypage}
\usepackage{setspace}
\onehalfspacing

% --- Typography ---
%\setmainfont{EB Garamond}           % replace with Latin Modern if unavailable
\usepackage[sfdefault]{sourcesanspro}
%\setmonofont{JuliaMono}[Scale=0.85]

% --- Mathematics ---
\usepackage{amsmath, amssymb, amsthm}
\usepackage{mathtools}

% --- Headers and footers ---
\usepackage{fancyhdr}
\pagestyle{fancy}
\fancyhf{}
\fancyhead[LE]{\small\itshape\leftmark}
\fancyhead[RO]{\small\itshape\rightmark}
\fancyfoot[C]{\small\thepage}
\renewcommand{\headrulewidth}{0.4pt}

% --- Chapter titles ---
\usepackage{titlesec}
\titleformat{\chapter}[display]
  {\normalfont\large\normalfont\centering}
  {\chaptertitlename\ \thechapter}
  {12pt}
  {\huge\centering}
\titlespacing*{\chapter}{0pt}{0pt}{30pt}

% --- Table of contents ---
\usepackage{tocloft}
\renewcommand{\cfttoctitlefont}{\normalfont\huge}
\renewcommand{\cftchapfont}{\bfseries}
\renewcommand{\cftsecfont}{\normalfont}
\setlength{\cftbeforechapskip}{6pt}


% --- Links and PDF ---
\usepackage[
  colorlinks=true,
  linkcolor=black,
  urlcolor=blue!60!black,
  citecolor=black,
  pdftitle={Whoever We May Be — The Projective Dynamic Logo},
  pdfauthor={Cédric Laubscher},
  pdfsubject={Scientific and philosophical popular science},
  pdflang={en}
]{hyperref}

% --- Epigraphs ---
\usepackage{epigraph}
\setlength{\epigraphwidth}{0.78\textwidth}
\renewcommand{\epigraphflush}{center}
\renewcommand{\epigraphrule}{0.4pt}

% --- Custom lists ---
\usepackage{enumitem}
\setlist[description]{leftmargin=2em, labelindent=0pt, font=\normalfont}

% --- Paragraph spacing ---
\setlength{\parindent}{1.4em}
\setlength{\parskip}{0pt}

% --- Custom commands ---
\newcommand{\ldp}{\textsc{pdl}}
\newcommand{\eg}{\textit{e.g.}\ }
\newcommand{\ie}{\textit{i.e.}\ }

% ============================================================
\begin{document}

% --- Title page ---
\begin{titlepage}
  \centering
  \vspace*{3cm}
  {\fontsize{11}{13}\selectfont Cédric Laubscher}\\[1.2cm]
  {\fontsize{30}{36}\selectfont\normalfont Whoever We May Be}\\[0.6cm]
  {\fontsize{16}{20}\selectfont\itshape
    An Introduction to the Projective Dynamic Logo}\\[0.4cm]
  {\fontsize{13}{16}\selectfont\itshape
    From the Electron to the Human Condition:\\
    the logic of coherence that governs everything}\\[3cm]
  \rule{0.45\textwidth}{0.4pt}\\[0.8cm]
  {\small Independent Researcher, Switzerland}\\[0.3cm]
  {\small March 2026}\\[1.5cm]
  \vfill
  {\footnotesize
    All technical publications of the PDL are available\\
    in open access on \href{https://zenodo.org}{Zenodo}
    under the name of Cédric Laubscher.}
\end{titlepage}

% --- Copyright page ---
\thispagestyle{empty}
\vspace*{\fill}
\begin{flushleft}
  \small
  \textcopyright\ 2026 Cédric Laubscher.\\[0.4em]
  This paper is made available under the
  \href{https://creativecommons.org/licenses/by/4.0/}{%
    Creative Commons Attribution 4.0 International (CC BY 4.0)} licence.\\[0.4em]
  You are free to share and adapt this text,
  provided you cite the author and the source.\\[1.2em]
  \textit{First edition, March 2026.}\\[0.4em]
  Typeset in \LaTeX\
\end{flushleft}
\vspace*{\fill}
\newpage

% --- Epigraph ---
\thispagestyle{empty}
\vspace*{4cm}
\epigraph{%
  Whoever we may be, from the atoms that form our molecules to our most intimate thoughts, we are nothing but the necessary expression of a simple and inexorable logic of coherence. A logic of causality that fulfils itself at every pulse. It is the Origin, the engine and the essence of the universe. The human being is merely its most vulnerable instrument.%
}{C.~Laubscher}
\newpage

% --- Table of contents ---
\frontmatter
\tableofcontents
\newpage

% ============================================================
\mainmatter
% ============================================================

% --- Foreword ---
\chapter*{Foreword \\ Why This Quest}
\addcontentsline{toc}{chapter}{Foreword}

There is a question I have never managed to put away in a drawer. Not the question of God, nor that of the meaning of life, although neither is far off. But a starker, more stubborn question: \emph{why does anything persist?}

\medskip

I am not talking about the cosmos in general, nor about life in particular. I am talking about the smallest and most ordinary thing there is: an electron. An electron exists. It has always existed, since the first seconds of the universe. It will still exist when our sun has gone out. It does not degrade. It does not age. It is rigorously identical to every other electron in the universe, at any place, at any moment. Why? Not \emph{how} --- physics is perfectly capable of describing how an electron behaves. But \emph{why} is it there? Why this particular form, and not another? Why this stability, which seems to defy entropy\footnote{A measure of the disorder of a physical system; according to the second law of thermodynamics, it can only increase in an isolated system.}?

\medskip

This is not only my question. It is the crack that physics itself has never managed to close: its two great theories, general relativity and quantum mechanics, contradict one another at the extremes --- at the heart of black holes, at the moment of the Big Bang. Our most powerful tools do not know what they are talking about where the question becomes truly bare.

\medskip

This question led me to do something perhaps a little mad: to start from scratch. Not to correct existing physics. To suspend it entirely. Remove space. Remove time. Remove particles. And ask the question in its absolute nakedness: under these conditions, what can \emph{exist}?

\medskip

The answer I found --- or rather that I watched emerge, because one does not choose these things --- is at once simple and vertiginous. It begins with the electron. It passes through the proton, through black holes, through the great constants that govern the universe. And it reaches all the way up to us.

\medskip

This paper is the account of that quest. Not a textbook. Not a scientific article. An account: the way in which a simple idea, a pulse, led step by step to a vision of the world that I invite you to explore together.

\chapter*{Abstract}
\addcontentsline{toc}{chapter}{Abstract}

\begin{quote}
\textit{Whatever we may be, from the atoms that form our molecules and our cells to our most intimate thoughts, we are nothing but the necessary expression of a simple and implacable logic of coherence. A logic of causality that fulfils itself at every pulse. It is the Origin, the engine and the essence of the universe. The human being is merely its most vulnerable instrument.}
\end{quote}

\bigskip

Contemporary physics lives with a contradiction it dare not name. Its two great theories --- general relativity, which governs the immensity of galaxies and black holes, and quantum mechanics, which governs the infinitely small world of electrons and protons --- are each of an almost miraculous precision. Yet they are fundamentally incompatible: their languages cannot coexist at the heart of a black hole, at the instant of the Big Bang, at the Planck scale. This is not a problem of calculation. It is a crack in the very foundations of our understanding of reality.

\medskip

The Projective Dynamic Logo (PDL) takes this crack seriously. It does not propose yet another correction. It does something more radical: it \emph{suspends everything}. No space. No time. No particles given in advance. And it poses the question in its absolute nakedness: \emph{in these conditions, what can exist?}

\medskip

The answer is at once simple and vertiginous. For something to exist in a lasting way~--- not just for an instant, but in a recognisable, real manner~--- it is necessary and sufficient that it form a cycle. A difference that returns upon itself. A pulse. That is all.

\medskip

From this single principle, without assuming either space or time, the PDL reconstructs step by step:

\begin{itemize}
\item The electron: the smallest structure capable of forming such a cycle~--- four entities, six links, one pulse~--- is precisely what nature calls an electron. Not by analogy: by logical necessity.
\item The proton: by stacking these cycles according to a precise hierarchy, five integers emerge, constrained and not adjusted, that describe the internal architecture of the proton with a rigour that standard physics has not yet achieved from its axioms.
\item The great constants of physics: $\hbar$, $\alpha$, $k_{\rm B}$, $G$ are no longer numbers fallen from the sky. They are \emph{coherence ratios}~--- the signatures of the way each structure manages its internal equilibrium. $\alpha$ and $\hbar$ are obtained without any freely adjusted parameter; $G$ requires the calibration of a single parameter $\varepsilon$, fixed once, then used to simultaneously recover $\alpha$ and $k_{\rm B}$.
\end{itemize}

One of the most striking images of the programme: an electron, a proton and a cosmic black hole all obey the same pulsative logic. In all three cases, what exists does so in the same way~--- a pulse, a closure upon itself, an active boundary that encodes everything the structure can say about itself. The electron is the minimal closure. The proton is its hierarchical version. The black hole is its extreme limit. These are not three objects of physics: they are three regimes of one and the same principle of existence.

\medskip

No complex structure can close perfectly upon itself~--- and it is precisely for this reason that the universe is granular. A cycle cannot be continuous if it must return \emph{exactly} to its initial state: discreteness is a logical constraint, not a hypothesis. There always remains a small leak of coherence: an incompressible fraction of openness towards what lies beyond the structure. This leak is the very mechanism by which gravitation emerges, and the structural condition of all transcendence. A cavity built solely from finite signed graphs spontaneously reproduces the density of states law $\rho(\nu)\propto\nu^{2}$, suggesting that the dimensionality of space could itself be a consequence of coherence constraints. A structural sketch further demonstrates that $(4,6)$ structures can provide a common language for general relativity and quantum mechanics~--- the two theories that physics has been unable to reconcile for a century. This result is preliminary, but it is structurally grounded: it does not postulate unification, it \emph{derives} it from the same principle of logical closure that selects the electron. If this result withstands future tests, the PDL will be the first programme to have reconstructed, from scratch and without a free parameter, a bridge between the infinitely small and the infinitely large.

\medskip

The human being~--- organism, identity, thought~--- does not escape this law: we are incomplete closures, kept alive by incessant exchanges with a world we do not entirely control. Our vulnerability is not an accident; it is the highest signature of the logic that governs the electron.

\medskip

This work is addressed to anyone who has wondered what it truly means to exist~--- and why this question concerns physics as much as philosophy, the black hole as much as the electron, the particle as much as oneself.

\newpage

% --- Introduction ---
\chapter*{Introduction \\ The Question Physics Dares Not Ask}
\addcontentsline{toc}{chapter}{Introduction}

There is a question that everyone has asked at least once, as a child or as an adult, gazing at the sky or staring into the void: \emph{why is there something rather than nothing?}

\medskip

Most of the time, one files this question away in the drawer of philosophy or religion, and moves on. Physics, for its part, seems to have decided not to concern itself with it: it describes \emph{how} things behave, not \emph{why} they exist. And yet something does not add up. Something profound, and which physics itself has not managed to resolve in a century.

\subsection*{A scandal that physics keeps quiet}

We have two great theories of the world. General relativity, formulated by Einstein in 1915, describes gravitation, the curvature of space-time, planetary orbits, the expansion of the universe, black holes. Quantum mechanics, developed in the 1920s, describes the behaviour of elementary particles, atoms, light, and all the physics of the infinitely small. Each of these theories is of a stupefying precision. Each has been verified millions of times, under ever more demanding experimental conditions. Each, taken on its own, is never truly wrong.

\medskip

And yet they are fundamentally incompatible. Not approximately, not in marginal cases: in their most fundamental concepts. General relativity presupposes a continuous, smooth, curved space-time; quantum mechanics presupposes discrete quantities, jumps, an irreducible indeterminacy. One describes a geometric and deterministic universe at large scales; the other a probabilistic and granular universe at small scales. These two languages cannot, in their current formulation, speak of the same reality.

\medskip

This is not a technical detail awaiting correction. It is an intellectual scandal that has lasted a hundred years. As soon as one pushes both theories towards their common limits~--- at the heart of a black hole, at the moment of the Big Bang, at the so-called Planck scale where gravitation and quantum effects simultaneously become unavoidable~--- their concepts collapse together. Space no longer makes sense. Time no longer makes sense. Particles no longer make sense. This is not a problem of calculation; it is a problem of \emph{language}. Our most powerful conceptual tools simply do not know what they are talking about at these extremes.

\medskip

Attempts at reconciliation are not lacking. String theory, loop quantum gravity, non-commutative geometry: immense programmes, mobilising hundreds of researchers for decades. None has yet produced a new experimentally verified prediction. None has resolved the underlying crack. This is not for want of intelligence or effort; it is perhaps because all these programmes correct existing theories without ever questioning their deepest shared assumption: that there exists, a priori, a space, a time, and entities that move within them.

\subsection*{Suspending everything}

The Projective Dynamic Logo (PDL) takes this problem seriously, and responds to it radically. It does not propose yet another correction to the existing equations. It does something more demanding: it \emph{suspends} everything. No space. No time. No particles given in advance. And it asks the bare question: under these conditions, what could still exist?

\medskip

The answer it proposes is at once simple and vertiginous. For something to truly exist~--- not merely for an instant, but in a lasting, recognisable, real way~--- it is necessary and sufficient that it form a cycle. A difference that returns to itself. A pulse. That is all.

\medskip

There is more. The same logic that gives rise to the electron also gives rise to the proton~--- and, beyond that, it sheds new light on one of the most fascinating objects in astrophysics: black holes. A black hole and a proton obey the same logic of surface: in both cases, physical information is encoded not in the volume but on a boundary. It is this structural kinship, derived and not postulated, that makes the PDL a serious candidate for the unification that physics has been seeking for a century. 

\medskip

One must measure what this means. For a century, the greatest minds in physics --- Einstein himself, then entire generations of theorists --- have sought to reconcile general relativity and quantum mechanics. String theory, loop quantum gravity, non-commutative geometry: immense programmes, without a single new experimentally verified prediction to date. The PDL does not arrive with yet another correction. It arrives with an answer to the prior question: \emph{why are there stable structures rather than nothing?} And it is precisely because it answers this question from first principles --- without presupposing space, time, or particles --- that the reconciliation between the two great theories is not an objective within it, but a \emph{consequence}.

\subsection*{A programme, not a doctrine}

The PDL does not stop at particle physics. It opens onto a broader vision: every thing that exists~--- an electron, a proton, a human being, a thought~--- is an incomplete closure. No structure can close perfectly upon itself. There always remains a leakage, an incompressible openness towards what lies beyond it. It is from this leakage that gravitation is born. It is this same openness that makes freedom, transcendence, and yes, our vulnerability possible.

\medskip

This paper is the first presentation of the PDL in English, written for the widest possible audience. You need no mathematical background to read it. Every technical concept will be translated into ordinary language before being used. Every analogy with human life will be offered as an invitation to reflect, never as a definitive certainty, but as an open proposition.

\medskip

The PDL is not a doctrine. It is a research programme: a rigorous and open way of asking questions that conventional physics sidesteps. Some of its results are solidly established. Others are precise conjectures, testable hypotheses. Others still are philosophical extrapolations, clearly identified as such. This honesty is, in our view, one of its greatest qualities.

\medskip

The PDL does not advance in isolation. It has engaged a formal dialogue with another fundamental research programme, the \emph{Theory of Objectivity} (TO), which seeks to identify the minimal logical conditions that any admissible universe must satisfy to allow the persistent existence of distinct objects. This dialogue~--- recorded in a dedicated publication within the corpus~--- has led to precise refinements: a more rigorous definition of active surfaces, a more explicit treatment of boundaries, and deeper reflection on what it means to exist for more than one observer. These exchanges do not constitute a merger of the two programmes; they illustrate that the PDL, far from being an isolated edifice, is sufficiently open and structured to engage with other fundamental approaches without losing its internal coherence.

\subsection*{What we shall build together}

We shall therefore begin at the beginning: nothingness, and the question of what can emerge from it. Not nothingness as a poetic backdrop~--- nothingness as a logical challenge. What does it mean to exist, truly, when one is no longer permitted to presuppose anything whatsoever?

\medskip

The answer, we shall build together, step by step. And at the end of the path, perhaps, the phrase that guides this entire text will take on a new meaning:

\begin{quote}
\itshape Whoever we may be, from the atoms that form our molecules to our most intimate thoughts, we are nothing but the necessary expression of a simple and inexorable logic of coherence. A logic of causality that fulfils itself at every pulse. It is the Origin, the engine and the essence of the universe. The human being is merely its most vulnerable instrument.
\end{quote}

% ============================================================
\chapter{Nothing, Then Something}

Imagine that you had to erase everything. Not just the objects around you, not just the Earth or the Sun or the galaxies, everything. Space itself. Time itself. Nothing remains, not even an empty background on which something could appear. Under these conditions, what could possibly happen?

\medskip

Most physical theories dodge this question. They assume that there already exists a space-time, already some laws, already potential particles; then they describe how all that evolves. This is legitimate, and it is remarkably effective. But it cheats with the original question. The PDL refuses this evasion. It deliberately places itself \emph{upstream} of all that and poses the problem in all its nakedness: what is necessary and sufficient for there to be something rather than nothing, in a lasting way?

\subsection*{Instantaneous existence does not count}

First, a brutal observation: a thing that exists only for an instant does \emph{not really count}. If a difference appears and then disappears without leaving the slightest trace, with no possibility whatsoever of being recognised or retrieved, then it is logically indistinguishable from nothingness. What has existed only once, without any possibility of returning, is as if it had never taken place.

\medskip

This is not an abstract philosophical statement; it is a logical constraint. For something to be real in the strong sense of the term, it must be \emph{identifiable}, that is to say retrievable, recognisable, re-instantiable. And this requires that it be capable of returning to itself.

\medskip

In the same way, a perfectly motionless state, absolutely uniform, without the slightest internal difference, is itself indistinguishable from nothingness: if strictly nothing happens, there is nothing to observe, nothing to distinguish, nothing to name.

\medskip

The only serious candidate for existence, in this austere framework, is therefore a difference that can repeat itself: a gap, a contrast, a distinction that is capable of returning to itself in a structured way. Not once~--- once is just chance. But a cycle, a return: that is something that can claim to exist.

\subsection*{The minimal pulse}

The PDL formalises this idea under the name of minimal binary pulse. Imagine two possible states: agreement and disagreement, yes and no, $+1$ and $-1$. A rule makes one pass from one to the other and back again: a perfect round trip, a cycle of length two. This is the absolute minimum of what one can call an event that lasts.

\medskip

Notice that, at this stage, there is still neither physical time, nor space, nor any measurable frequency. There is only the logical possibility of an alternation: something can return to what it was. That is all. But it is already infinitely more than nothing. I remember the moment when this idea struck me for the first time. Not in a laboratory, not in front of an equation, but in the silence of an ordinary night, wondering why an electron does not age. The answer I eventually glimpsed was so simple that it first seemed derisory. Then vertiginous.

\medskip

One can find an image in everyday life. A heart that beats does not exist because of the space in which it is located; it exists because it can repeat its cycle indefinitely. A melody does not exist because it occupies a place; it exists because one can recognise it, play it again, find it again. The PDL says: at bottom, everything that exists operates according to this same principle. Space and time come \emph{afterwards}, as tools for organising and comparing cycles, not before, as a pre-existing backdrop.

\subsection*{A universe granular by necessity}

There is a consequence that one does not always notice at first glance, but which is fundamental. If existence rests on a cycle~--- on an exact return to the initial state~--- then that cycle cannot be continuous.

\medskip

Here is why. A continuous process, in the mathematical sense, can always be interrupted halfway: it passes through an infinity of intermediate states, and nothing guarantees that it will ever find its starting point again exactly. A cycle that must close \emph{exactly}, without drift, without accumulation of error, must on the contrary be discrete: it passes through a finite number of well-distinct states, and returns to one of them rigorously. That is the only way for a structure to remain identical to itself through time.

\medskip

In other words: the granularity of the universe is not a hypothesis that the PDL adds to its model. It is a logical constraint that the programme \emph{derives}. If something must exist in a lasting way, it must be discrete. Quantum physics has accustomed us to the idea that energy is transmitted in packets, that the energy levels of an atom are discrete, that nature does not know the infinitely divisible at every scale. The PDL says this is no accident: it is the inevitable signature of everything that can maintain itself.

\medskip

This idea resonates with the deepest currents of contemporary physics~--- holography, entropic gravity, quantum information~--- which all converge on the same observation: the continuous is a convenient approximation, and deep reality is combinatorial. The PDL situates itself within this current, and offers a justification from first principles.

\subsection*{Why a single pulse is not enough}

An isolated pulse, however, poses a problem. If it is connected to nothing else, how can it be distinguished from another identical pulse? How can it \emph{refer} to itself? Without a relational context, it remains undefinable: two identical and isolated alternations are strictly indistinguishable from one another.

\medskip

For a distinction to be operational, it must be anchored in a network of relations, it must be able to situate itself with respect to something else. This is why the PDL introduces a second requirement: elementary entities must be linked to one another, and these relations must themselves form closed loops.

\medskip

The simplest closed loop is a triangle: three entities, each linked to the other two. A triangle is the smallest circuit capable of closing upon itself without appealing to anything external. And in a triangle, something remarkable can happen: the three relations can be mutually \emph{coherent} or \emph{incoherent}, depending on how their signs combine.

\medskip

The PDL then posits a central axiom: admissible structures are those that minimise triangular incoherence. A structure that fails to satisfy its own internal constraints cannot maintain itself; it dissolves. Only those configurations persist that manage to close upon themselves in a coherent way.

\subsection*{The first admissible structure}

We can now pose the concrete question: what is the smallest network of linked entities that simultaneously satisfies the two requirements~--- pulse and triangular coherence~--- without appealing to any external support whatsoever?

\medskip

The answer is unique, and its precision is surprising: four entities linked by six edges. Neither fewer, nor more. With three entities, one can form a coherent triangle, but one cannot propagate reference across a sufficiently rich network: a single loop is too fragile. With four entities and six edges, a logical tetrahedron, denoted $(4,6)$, one reaches the first threshold at which everything fits together: all loops are coherent, the pulse cycle is possible, and the structure is self-sufficient.

\medskip

This $(4,6)$ structure is not invented or chosen arbitrarily: it is \emph{derived} from the constraints. It is the first logically obligatory answer to the question: what can exist without presupposing anything else?

\medskip

In the next chapter, we shall see that this abstract structure~--- four points, six links, one pulse~--- is precisely what the PDL identifies as the logical archetype of the electron. Not because the electron \emph{is} a mathematical graph, but because the electron is, in nature, the first physical realisation of what the logic of coherence makes obligatory.

\medskip

Before we get there, let us pause for a moment on what we have just done. We have used neither space, nor time, nor energy, nor any law of physics. We have only asked: what can return to itself in a coherent way? And a precise structure has emerged, inevitable. That is the heart of the PDL: existence is not a backdrop. It is a problem of logical optimisation, and physical reality is its solution.

% ============================================================
\chapter[The First Pulse --- The Electron]{The First Pulse \\ The Electron}

We established in the previous chapter that the smallest structure capable of existing in an autonomous and coherent way is the (4,6) network: four entities, six links, one pulse cycle. This structure was not invented; it was derived from the most minimal constraints one can impose on anything that claims to exist durably.

\medskip

But one question arises immediately: does this purely logical construction have anything to do with the real world? The PDL’s answer is yes, and in a very precise way.

\subsection*{The electron, discreet star of the universe}

Among all known particles, the electron occupies a special place. It is strictly stable: unlike many particles that decay within fractions of a second, an isolated electron never transforms into anything else. It is universal: all the electrons in the universe are rigorously identical, at any place, at any moment. It is structurally simple: at all scales accessible to our instruments, it has no detectable substructure. And it is animated by an internal oscillation, the Compton frequency\footnote{a frequency related to the mass of the electron by $f = mc^2/h$; it is the natural rhythm at which the electron “beats” in its own rest frame.}, which is its own and is directly linked to its mass.

\medskip

These four characteristics — stability, universality, absence of internal structure, intrinsic oscillation — are precisely those one expects from the (4,6) structure. Stable because its internal coherence is maximal and does not degrade. Universal because it is the only solution to the logical problem posed. Without internal structure because it is minimal by construction. And animated by a pulse because that is its very definition.

\medskip

The PDL therefore proposes what it calls a minimal correspondence principle: the electron at rest is the first physical realisation of the logical closure (4,6). This is not a metaphysical identity. One does not say that the electron \emph{is} a graph. It is a precise and testable structural hypothesis.

\medskip

To understand it, one must know what is meant by a \emph{stationary regime}: a structure is said to be stationary when it beats, spins, oscillates, but without ever transforming into anything else. It changes state at each pulse, but it always returns to the same initial pattern. It is like a spinning top: it is in permanent motion, but its overall shape remains stable. By contrast, a \emph{non}-stationary structure degrades, disintegrates, or transforms into something else over time.

\medskip

In this framework, the question becomes: if something must exist in nature as the \emph{first} stationary regime, that is, the simplest structure that beats without ever degrading, what is the natural candidate? The PDL’s answer is the electron: stable, universal, without substructure, animated by a permanent internal oscillation. Everything one expects from the (4,6) closure, dressed in physics.

\subsection*{The cost of existing}

This correspondence is not just a pleasing analogy. It allows us to reread an equation everyone has seen without ever truly understanding it:

\begin{equation}
E = \hbar\,\omega
\end{equation}

Before reading it, a word on the term \emph{quantum} — plural: \emph{quanta}. In everyday life, energy seems able to take any value: one can heat water a little, a lot, very slightly. But at the scale of particles, nature does not work that way. Energy is not continuous like flowing water; it is transmitted in discrete packets, minimal portions that cannot be subdivided further. A quantum is precisely such a packet: the smallest amount of energy that a given structure can exchange in a single go. One cannot receive half a quantum, any more than one can receive half a coin in a transaction that demands it whole.

\medskip

In ordinary physics, the equation $E = \hbar\,\omega$ simply states that the energy of this packet is proportional to the frequency $\omega$ of the associated oscillation, with $\hbar$ (the reduced Planck constant, pronounced “h-bar”) as the proportionality factor. This is correct, but it is mute: why this relation? Where does $\hbar$ come from? Why this particular value?

\medskip

In the PDL, the same equation is read differently. $\omega$ is the pulse frequency of the (4,6) structure, the physical expression of its logical pulsation. And $\hbar$ is no longer a mysterious parameter imported from outside: it is the minimal cost of a coherence cycle, the amount of action required for the structure to loop back onto itself once. To exist, even in the simplest possible way, has a price. This price is paid at each pulse. And $\hbar$ is the currency.

\medskip

This rereading transforms what seemed to be an arbitrary law into a logical necessity. The Planck constant is not a curiosity of nature; it is the quantitative translation of the fact that no existence cycle is free.

\subsection*{Mass as a coherence budget}

If $\hbar$ is the cost of a cycle, then mass and energy can be reread as coherence budgets: measures of how many cycles a structure must maintain, and at what rate, in order to continue to exist.

\medskip

For the simple (4,6) structure, the electron, this budget is minimal and uniform. All its coherence is internal, well organised, without hierarchy. This is why the electron is so light: it pays the bare minimum to exist.

\medskip

Inertia, too, finds a new reading here. Why is it difficult to change the state of motion of a massive object? Because such an object is a structure whose internal coherence is dense and strongly organised. Modifying it requires reorganising a large number of cycles simultaneously, which demands a proportional investment of coherence. What we call mass is in reality the \emph{internal rigidity} of a logical closure.

\subsection*{A structure, not a point}

One of the most persistent habits of classical physics is to treat elementary particles as \emph{points}: entities without dimension, without interior, without structure. The PDL breaks with this habit. Not by assigning them a geometric size, but by recognising in them an internal relational organisation.

\medskip

The electron, in this framework, is not a point in space. It is a pattern of relations, four nodes, six links, a cycle that sustains itself. What matters is not where it is, but \emph{how it is organised}. The space in which it is located is an emergent description, convenient for comparing several electrons with one another, but no more fundamental than the relational structure that defines it.

\subsection*{The first instrument}

Let us return to the sentence that guides this paper. It says that we all are, from atoms to thoughts, instruments of a logic of coherence. The electron is the first of these instruments: the simplest way the universe has of concretely realising what logic makes obligatory.

\medskip

It is not there by chance. It is not there because an external law created it. It is there because, as soon as something must exist in a minimal, autonomous and coherent way, the (4,6) form is inevitable, and the electron is that form, dressed in physics.

\medskip

But the universe does not content itself with electrons. It builds far more complex structures: protons and neutrons, nuclei, atoms, molecules, stars, living beings. How does one pass from the minimal pulse to all this richness? That is the question of the next chapter.

\medskip

But the electron is not only defined by its structure; it is also defined by the way it responds when one interrogates it. And here again, the logic of coherence imposes a unique answer.

\subsection*{What measurement owes to coherence}

There is a question that quantum mechanics has been asking for a century without ever truly answering: why are quantum probabilities \emph{squares}? Why, when one measures the spin\footnote{an intrinsic property of quantum particles, analogous to a rotation on itself but with no classical equivalent; it takes only two discrete values, “up” or “down”.} of an electron, is the probability of obtaining a result proportional to the square of an amplitude, and not to the amplitude itself, nor to its cube?

\medskip

This rule, called Born’s rule, after the physicist Max Born who formulated it in 1926\footnote{%
    Formally, if a quantum system is in the state $|\psi\rangle = a\,|\!\uparrow\rangle + b\,|\!\downarrow\rangle$, where $a$ and $b$ are complex numbers satisfying $|a|^2 + |b|^2 = 1$, then the probability of obtaining the result “spin up” in a measurement is $|a|^2$, and that of obtaining “spin down” is $|b|^2$. It is the square of the modulus of the amplitude — and not the amplitude itself — that gives the probability. This rule applies to all observables in quantum mechanics, not just spin: the probability of finding a particle at a point in space, the probability that an atom emits a photon, the probability that a nuclear reaction occurs — all this is calculated via squared amplitudes. Why squares and not something else? In standard physics, this rule is simply postulated; the PDL proposes, for the first time in a discrete framework, a structural derivation of it.} 
is one of the cornerstones of all quantum physics. It is confirmed with extraordinary precision in millions of experiments. But in the standard framework, it is \emph{postulated}: one lays it down as a rule of the game, without explaining its origin. Why squares? Nobody truly knows. It is one of the questions that resisted me the longest. I circled around it for months before understanding that the answer was not in the existing equations, but upstream of them, in the very structure of what it means to measure something that beats.

\medskip

That is where the opening came from. The PDL gives an answer, and it is one of the programme’s most solid pieces. Within it, measuring the spin of an electron means coupling the (4,6) structure to a measuring device, itself modelled as a network of relations. This coupling is made through what are called mixed triangles: triplets of links that connect two nodes of the electron to one node of the device. These triangles can be coherent or incoherent, exactly like the internal triangles of the (4,6) structure. And the same principle applies: the system globally minimises incoherence.

\medskip

The device has two branches, say “spin up” and “spin down”. Each configuration of the global system is weighted according to the number of coherent mixed triangles it realises in each branch. In the limit where this optimisation is strong, the effective probabilities of ending up in one branch or the other converge towards a precise expression. And this expression is exactly Born’s rule: the probabilities are the squares of the amplitudes.

\medskip

This result was then consolidated by a Gleason-type proof\footnote{Gleason's theorem, established by the American mathematician Andrew Gleason in 1957, answers the following question: if one wants to assign probabilities to the results of quantum measurements in a coherent way~--- that is, in such a way that the probabilities do not depend on the measurement context chosen~--- what form can these probabilities take? The answer is remarkably constraining: in a Hilbert space of dimension at least three, the only possible assignment is precisely Born's rule. Not one among others: the only one. This result is regarded as one of the deepest demonstrations in the foundations of quantum mechanics, because it shows that Born's rule is not an arbitrary choice~--- it is the inevitable consequence of demanding minimal coherence in the way probabilities are assigned. The PDL establishes an analogous result in a discrete, combinatorial framework rather than in the continuous Hilbert space (an abstract mathematical space in which quantum states are represented as vectors; the standard framework of quantum mechanics) of standard physics: this is what makes it a structural advance, and not a mere reformulation (\textit{Toward a Gleason-Type Uniqueness Theorem for Born's Rule in the PDL Framework}, C.~Laubscher, Zenodo, 2026).}. Concretely, under natural operational conditions~--- phase invariance, rotational covariance, repeatability of measurement, symmetry between the two branches~--- it is possible to show that Born's rule is the only assignment of probabilities compatible with the PDL structure. It is not merely that it emerges naturally: no other rule can exist in this framework without contradicting its own constraints.

\medskip

In other words: in a universe governed by the PDL's logic of coherence, quantum probabilities \emph{could only be squares}. This is not a choice. It is a structural necessity, and its consequences are far-reaching: the most fundamental rule of quantum mechanics~--- the rule that tells how the probabilistic world of the infinitely small translates into observable outcomes~--- is not an arbitrary convention. It is the inevitable consequence of a universe that optimises its relational coherence. The same logic that selects the electron as the first admissible structure also dictates the way it responds to a measurement.

% ============================================================
\chapter[Building Bigger --- The Proton]{Building Bigger \\ The Proton}

The electron is admirable in its simplicity. But the universe we inhabit does not reduce to isolated electrons. There are nuclei, atoms, molecules, and at the heart of all this there is the proton, about two thousand times heavier than the electron and with an internal complexity of a radically different kind.

\medskip

How does the PDL approach the proton? Not by postulating it, not by adding it in by hand. By \emph{deriving} it, by showing that, once one stacks $(4,6)$ closures according to the same criteria of coherence and optimisation, a precise hierarchical architecture imposes itself.

\subsection*{One cannot just add electrons}

The first temptation would be to say: the proton is simply several $(4,6)$ structures placed end to end. But that does not work. If each substructure seeks to maximise its own internal closure, the whole becomes too dense, too saturated. There is no longer room for a rich internal dynamics, no more flexibility, no more possibility of interaction with the outside. A sum of absolutes does not make a coherent whole.

\medskip

The PDL therefore imposes an additional rule: in a composite structure, there must be a hierarchy. Some regions are strongly closed; these are the stationary cores. Other regions are more open, more dynamic; this is the relational sea. This asymmetry between local closure and global openness is what allows the structure to remain coherent without freezing.\footnote{Readers familiar with particle physics will recognise structural analogues here: the stationary cores correspond to \emph{valence quarks}, and the relational sea to the field of \emph{gluons} and \emph{virtual quarks} of quantum chromodynamics. But these technical terms are not necessary to follow the argument: the PDL reconstructs this architecture independently, without postulating it.}

\subsection*{Integers that are not parameters}

When the PDL applies its optimisation criteria to this architecture, namely: maximal coherence, minimal hierarchy, net charge compatible with that of the electron, it obtains a structure characterised by five precise integers:

\begin{itemize}
\item $n_u = 24$ and $n_d = 28$: the sizes of the two types of stationary cores (up and down quarks), expressed in number of elementary entities.
\item $r_{\text{val}} = 930$: the total number of internal links in the three valence cores.
\item $R_{\text{mer}} = 10\,087$: the number of links in the dynamic relational sea.
\item $r_{\text{total}} = 11\,017$: the total relational budget of the proton.
\end{itemize}

These numbers are not tuned to match experimental data. They are \emph{selected} by logical constraints: divisibility by four, controlled asymmetry between the cores, and maximisation of global coherence under a minimality constraint. The closeness to experimental results, in particular the $9\,\%$ stationary / $91\,\%$\footnote{The approximate 9\,\% stationary / 91\,\% dynamic split is consistent with the results of deep inelastic scattering (DIS) experiments. Combined measurements from the H1 and ZEUS experiments at the HERA collider have established that valence quarks carry only about 45\,\% of the proton's total momentum, the rest being attributed to gluons and sea quarks. See: H1 and ZEUS Collaborations, \textit{Combined Measurement and QCD Analysis of the Inclusive $e^{\pm}p$ Scattering Cross Sections at HERA}, \textit{JHEP} 01 (2010) 109, arXiv:0911.0884 [hep-ex].} dynamic partition of the proton's energy, is an \emph{a posteriori} validation, not an adjustment.

\medskip

These five integers did not fall from the sky, nor were they adjusted by hand to fit the data. They are selected by the PDL optimisation functional under strictly combinatorial constraints\footnote{The derivation of these five integers, the systematic numerical simulations scanning the space of neighbouring configurations, and the local uniqueness conjecture are developed in full in \textit{On the Combinatorial Selection of the Proton Architecture in PDL --- Axiomatic Constraints, Integer Structures, and a Uniqueness Conjecture}, C.\,Laubscher, Zenodo, 2026. A formal proof of global uniqueness remains an open problem.}: divisibility by four, maximal triangular coherence, stationary / dynamic fraction close to $9/91$, net charge $+1$, and simultaneous compatibility with the fine-structure constant. Systematic numerical simulations have been carried out\footnote{These explorations are the subject of a dedicated article in the PDL corpus: \emph{On the Combinatorial Selection of the Proton Architecture in PDL --- Axiomatic Constraints, Integer Structures, and a Uniqueness Conjecture} (C.~Laubscher, Zenodo, 2026). The simulations, performed on Google Colab, scanned the space of neighbouring configurations by varying the multiplicities $n_u$ and $n_d$ in steps of four and adjusting $R_\text{mer}$ in moderate intervals. In each case, the alternative configurations violated at least one constraint among the $9/91$ partition, the accuracy on $\alpha^{-1}$, or global triangular coherence. These results provide strong numerical support for the local uniqueness conjecture, without constituting a formal proof: a uniqueness theorem in the strict mathematical sense remains an open problem.} to test the rigidity of this quintuplet. Changing a single integer, for example taking $r_\text{total}$ from $11\,017$ to $11\,018$, is enough to destroy the simultaneous agreement with the constants of nature. This is not a numerical coincidence: it is the signature of an extremely tight structural constraint. The proton architecture is therefore the most constrained candidate the programme has identified to date, but the formal and exhaustive demonstration of its absolute uniqueness remains an open problem and one of the most stimulating mathematical challenges that the PDL poses to specialists in combinatorial optimisation.

\medskip

I will not forget the moment when these five numbers ceased to be variables and became constraints. When one realises that one cannot change a single one without destroying everything, something changes in the way one looks at physics, and perhaps at reality as a whole.

\subsection*{The active surface and the golden ratio}

The proton does not live in isolation. It interacts with electrons, with other protons, and with radiation. For this it has an active surface: a subset of its relations projected outwards and made available for couplings with other structures.

\medskip

The solution to the optimisation problem for this surface brings out, in a conjectural but strongly constrained way, the golden ratio\footnote{A ratio that appears naturally in problems of self-similarity and optimal partition.} $\varphi = (1+\sqrt{5})/2$. Not as a mystical symbol, but as the value that optimises the partition between inside and outside in a hierarchical structure subject to the PDL constraints.

\subsection*{The constants of nature as coherence ratios}

It is at this stage that the PDL accomplishes something remarkable. The great physical constants here find a precise structural interpretation.

\medskip

\medskip

The fine-structure constant $\alpha \approx 1/137$ is reinterpreted as the ratio between the active surface of the proton and its total relational budget: the fraction of coherence that the proton projects outward in order to interact with an electron.\footnote{$\alpha$ (alpha) is one of the most famous constants in physics. It measures how strongly charged particles~--- such as the electron and the proton~--- attract or repel one another via the electromagnetic force. Its value, about $1/137$, is a dimensionless number: it does not depend on any system of units. Why precisely $1/137$ and not something else? Nobody knows within standard physics. This is precisely what the PDL seeks to explain.} Boltzmann's constant $k_B$ encodes the way coherence is distributed between microscopic and macroscopic levels of organisation.\footnote{$k_B$ (pronounced "Boltzmann's constant", after the Austrian physicist Ludwig Boltzmann, 1844--1906) is the bridge between two worlds: that of individual particles and that of the heat we feel. It translates what we call \emph{temperature}~--- a macroscopic notion, that of the thermometer~--- into the average agitation of particles at the microscopic scale. Without it, statistical physics and thermodynamics could not communicate.} The gravitational constant $G$, finally, is obtained by hierarchical amplification of this same leakage rate $\varepsilon$ across eighteen levels of organisation.\footnote{The detail of these eighteen filters, organised in three blocks of six levels (proton interior, proton-to-nucleus, nucleus-to-macroscopic matter), together with the derivation of the effective gravitational coupling, is developed in \textit{Coherence Leakage, Hierarchical Filtering, and the Exponent\,18 in the PDL Framework}, C.\,Laubscher, Zenodo, 2026.} It is important to be precise about the status of this result: $G$ is not derived without a free parameter as $\alpha$ is. The value of $\varepsilon$ is, at this stage of the programme, adjusted to reproduce the experimental value of $G$; it is then this same calibrated $\varepsilon$ that is used to recover $\alpha$ and $k_B$ simultaneously. What the PDL establishes structurally is the \emph{form} of the link between $G$, $\alpha$ and the proton architecture; the numerical \emph{value} of $\varepsilon$ remains, for now, an empirical parameter. This distinction between a fixed logical architecture and partial numerical calibration is one of the marks of epistemic honesty of the programme.

\medskip

What is essential to retain: the laws of nature are not arbitrary rules inscribed in the cosmos. They are compact descriptions of the way a finite universe manages its coherence budgets at all scales, from the electron to the stars.

\subsection*{The proton and the black hole: a shared logic of surface}

There is a kinship that standard physics has not yet fully recognised, and which the PDL brings to light naturally: the proton and the black hole obey the same fundamental principle. In both cases, what matters physically is not encoded in the volume, but on a boundary.

\medskip

For the black hole, this principle is known as holography. Since the work of Bekenstein and Hawking in the 1970s, it has been established that the entropy of a black hole is proportional to the area of its horizon, not to its volume. In other words, all the physical information contained in a black hole is inscribed on its surface, like a hologram. This result, confirmed by numerous independent approaches, suggests that the dimensionality of the interior volume is an emergent description, not a fundamental one.

\medskip

The PDL recovers exactly the same principle for the proton, but from the logic of coherence: the active surface of the proton, that fraction of its relations projected outwards, is the only place where the proton can be interrogated, measured, or coupled to anything else. Its interior~--- the five integers, the stationary cores, the relational sea~--- is structurally inaccessible from outside except through what its surface projects. This is not an instrumental limitation: it is a logical necessity. A closure can interact with what lies beyond it only through what it exposes.

\medskip

The difference between the two objects is one of scale and direction: the black hole draws macroscopic coherence inward towards its horizon, while the proton projects its coherence outward through its active surface. One is a closure that absorbs; the other a closure that radiates. But the structure is the same: information on the boundary, inaccessible interior, and an incompressible leakage that governs interaction with the rest of the world.

\medskip

This convergence is not a metaphor, and it is not postulated. It emerges from the same principle of logical closure applied to two radically different regimes of organisation. If it withstands testing, it constitutes one of the strongest arguments for a structural unification of particle physics and gravitation~--- not by adjoining equations, but by derivation from a common principle.

\subsection*{An aside: what the PDL says about computation}

Readers familiar with computer science may have noticed a troubling resonance. The $(4,6)$ structure~--- a network of four nodes linked by six edges subject to a discrete update rule~--- looks remarkably like a minimal cellular automaton. This is no accident. The PDL suggests that computation and physical existence share a deep common structure: in both cases, what runs stably is what \emph{exists}. A programme that loops indefinitely without halting is, in this framework, the software version of a closure; a programme that crashes is a structure that has been unable to maintain its internal coherence.

\medskip

This resonance is not just a metaphor. Recent work on quantum complexity and quantum circuits shows that the stability of a quantum computation depends precisely on the ability of logic gates to maintain entanglement against environmental decoherence. The central problem of quantum computing~--- how to preserve coherence long enough to perform a useful computation~--- is, in the language of the PDL, the problem of maintaining an open closure without dissolving it. Quantum error-correcting codes, which detect and repair coherence violations without directly measuring the state of the system, can be read as technological implementations of the minimal-leakage principle.

\medskip

\medskip

More generally, the PDL invites a reformulation of the fundamental question of computational complexity: why are some problems harder than others? In the relational framework of the PDL, the difficulty of a computation could be related to the coherence cost required to maintain all the constraints of the problem simultaneously. The complexity classes $\mathbf{P}$ and $\mathbf{NP}$ could then reflect a structural difference in the coherence budgets required~--- an open but precise conjecture.

\subsection*{A universe that builds itself}

The proton marks a decisive threshold. Beyond the first minimal closure that is the electron, here is a composite architecture that emerges from a hierarchy of organised pulses: stationary cores, a dynamic sea, an interaction surface, constants that are its imprints.

\medskip

But there is a limit to this nesting logic. No structure, however well organised, can close completely upon itself. There is always something that escapes~--- an incompressible leakage, an irreducible opening towards what lies beyond it. It is from this leakage, precisely, that gravitation is born. And it is this, too, that makes us vulnerable.

% ============================================================
\chapter[Nothing Ever Closes Completely --- Leakage]{Nothing Ever Closes Completely \\ Leakage}

Up to this point, the PDL has built: first the minimal pulse, then the electron, then the proton. At each stage, the logic of coherence selects the most closed, the most stable, the most autonomous structure possible. One might think that the ultimate goal is perfect closure, a structure that would need nothing, that would maintain itself indefinitely without ever letting anything leak out.

\medskip

This perfect closure does not exist. And this is not a defect of the programme; it is one of its deepest results.

\subsection*{Irreducible incoherence}

Let us recall the principle of triangular coherence: in any admissible structure, triangles of relations must satisfy a condition of internal agreement. For the (4,6) structure of the electron, this condition can be satisfied perfectly: all triangles are coherent, no internal contradiction remains. This is why the electron is so stable.

\medskip

But as soon as one builds a composite structure, as soon as one assembles several elementary closures into a richer architecture, something changes. The number of triangular constraints to be satisfied simultaneously grows much faster than the number of available degrees of freedom. Beyond a certain level of complexity, it becomes \emph{mathematically impossible} to satisfy all constraints at once.

\medskip

Imagine a group of five friends who all want to get on perfectly with one another. Between two people, agreement is possible; one only has to find common ground. But between five people, there are ten pairs of relations to manage simultaneously. And experience shows that the larger the group becomes, the harder it is for everyone to get along with everyone else at the same time: some tensions become inevitable, not through ill will, but because the requirements of each relationship end up contradicting one another. One can minimise frictions, distribute them intelligently, but one cannot eliminate them all. This is exactly what the PDL says for physical structures: beyond a certain threshold of complexity, a residue of incoherence is structurally inevitable. This residue is coherence leakage.

\medskip

There is always a residue left: a fraction of triangles that cannot be coherent, whatever one does. The PDL quantifies this fraction by a parameter denoted $\varepsilon$, the coherence-leakage rate. For the proton, this rate is very small, but it is strictly positive. It cannot be brought to zero. It is structural.

\subsection*{What leakage becomes}

A first reaction would be to see $\varepsilon$ as a mere flaw, a residual imperfection without major consequence. The PDL says the opposite: this leakage is not lost. It is \emph{exported}.

\medskip

The coherence that cannot be absorbed inside a structure propagates outward, into the network of relations that surrounds that structure. It does not disappear; it becomes a constraint on neighbouring structures. On a large scale, this exported coherence manifests itself as a modification of what is called the metric, the way in which structures influence one another at a distance.

\medskip

In other words: gravitation is the macroscopic manifestation of coherence leakage\footnote{\textit{Resonance~--- Bekenstein, Hawking, and surface entropy (1972--1974). In 1972, Jacob Bekenstein showed that the entropy of a black hole is proportional to the area of its horizon~--- not its volume. Stephen Hawking completed this in 1974 by showing that black holes emit thermal radiation at a temperature that links simultaneously $G$, $\hbar$, and $k_B$~--- the three constants the PDL connects through a single leakage parameter~$\varepsilon$. This is not yet a derivation: the PDL does not yet explain Hawking radiation. It is a structural resonance: the idea that physical information resides on a boundary rather than in a volume is precisely what the proton's active surface instantiates in the PDL\@. A formal dialogue between these two frameworks remains to be built.}}. 
What Newton described as a force of attraction between masses, what Einstein reinterpreted as a curvature of space-time, the PDL reads as the collective propagation of the coherence irreducibly lost by each composite structure.\footnote{%
  This claim does not rest on intuition or analogy: it is the result of a detailed mathematical derivation, developed in several publications of the PDL programme (\textit{Coherence Leakage, Hierarchical Filtering, and the Exponent~18 in the PDL Framework}; \textit{A Structural Derivation of the Fine-Structure Constant from the PDL}; \textit{From Discrete Coherence Flux to Effective Fields}, all available on Zenodo). What gives it force, and what distinguishes it from a simple narrative reformulation, is the nature of the convergence obtained. The same formal framework, starting from the same logical constraints and without any \emph{ad hoc} fitted parameter, derives simultaneously and independently three quantities:
  \begin{itemize}
  \item the coherence-leakage rate $\varepsilon$, quantified at the scale of the proton;
  \item the fine-structure constant $\alpha \approx 1/137$, expressed as a combinatorial ratio on the proton structure;
  \item the gravitational constant $G$, obtained by hierarchical amplification of $\varepsilon$ across about eighteen levels of organisation.
  \end{itemize}
\textit{Epistemic note:} among these three quantities, $\alpha$ is derived without any freely adjusted parameter; the value of $\varepsilon$ is, by contrast, calibrated to reproduce $G$, and this same $\varepsilon$ is then used to test coherence with $\alpha$ and $k_B$. The strength of the result lies in this \emph{mutual coherence}: a single parameter, fixed once, simultaneously reproduces three constants that standard physics treats as entirely independent. These three values are compatible with experimental measurements. Their mutual coherence — the fact that they emerge from the same logical edifice without contradicting one another — constitutes the true keystone of the programme. In physics, recovering one fundamental constant by derivation is already remarkable. Recovering three simultaneously, from a single principle of coherence and without freedom of adjustment, is the most stringent test one can impose on a theory — and the PDL passes it.}

\subsection*{Hierarchical filtering and the exponent 18}

How can leakage so tiny at the scale of the proton produce an effect as observable as gravitation?

The PDL’s answer is a mechanism of \emph{hierarchical filtering}: leakage does not produce gravitation directly. It produces it after passing through eighteen successive levels of organisation. At each level, leakage from the lower level is slightly amplified and transmitted to the higher level. After eighteen steps, the signal, starting from almost nothing at the scale of a proton, has acquired the intensity we measure as gravitation.

Here, concretely, are the main stages of this cascade\footnote{%
  The detail of these eighteen filters, organised in three blocks of six, is developed in \emph{Coherence Leakage, Hierarchical Filtering, and the Exponent~18 in the PDL Framework} (C.~Laubscher, Zenodo, 2026).}:

\begin{itemize}
  \item Levels 1–6: the interior of the proton. Leakage is born at the heart of the proton architecture itself — between the valence cores, the dynamic sea, and the active surface. It is here that the parameter $\varepsilon$ is defined and quantified.
  \item Levels 7–12: from proton to atomic nucleus. Protons and neutrons assemble into nuclei. The leakage of each nucleon combines and is redistributed at nuclear level. Nuclear stability constraints constitute the filters of this block.
  \item Levels 13–18: from nucleus to macroscopic matter. Nuclei form atoms, atoms form molecules, molecules form macroscopic objects, then self-gravitating structures. At each step, accumulated leakage is amplified according to the laws of collective coherence until, at the end of the process, it produces an effective coupling of order $10^{-38}$ — that is, precisely what we measure for the gravitational constant $G$.
\end{itemize}

\medskip

This result deserves a pause. Conventional physics treats $\alpha$ and $G$ as two entirely independent constants\footnote{In 2025, a study published in \textit{AIP Advances} (Vopson \textit{et al.}) independently proposed that gravitation could emerge from an entropy-reduction process in a computational universe --- a striking conceptual convergence with the PDL's coherence-optimisation principle, reached by an entirely different route.}: no one has ever known why they take the values they do, nor why they should have the slightest link between them. The PDL says the opposite: these two constants are two faces of a single combinatorial defect in the proton architecture. There is only one leakage, only one parameter $\varepsilon$; depending on whether one projects it towards the active surface (electromagnetic coupling) or towards hierarchical filtering (gravitation), one obtains $\alpha$ or $G$\footnote{Resonance~--- Verlinde and entropic gravity (2010) In 2010, Dutch physicist Erik Verlinde proposed that gravitation is not a fundamental force but emerges thermodynamically from entropy variations --- exactly as gas pressure emerges from the agitation of its molecules(E.\ Verlinde, \textit{On the Origin of Gravity and the Laws of Newton}, JHEP \textbf{04} (2011) 029, arXiv:1001.0785; M.\,M.\ Brouwer \textit{et al.}, \textit{First test of Verlinde's theory of emergent gravity using weak gravitational lensing measurements}, MNRAS 466, 2547 (2017)). Observational tests on more than 33{,}000 galaxies in 2016--2017 confirmed partial agreement with this prediction, without requiring dark matter. The PDL mechanism is structurally more precise: where Verlinde postulates an elastic volume entropy, the PDL proposes an explicit combinatorial derivation of the leakage $\varepsilon$, amplified across eighteen hierarchical levels. Both approaches converge towards the same central intuition --- gravitation as \emph{bookkeeping of information on boundaries} --- but by entirely independent paths.}.

\medskip

In other words, what is remarkable is not only the final result. It is that the same parameter $\varepsilon$, defined once and for all at the scale of the proton, simultaneously and without free parameter reproduces the value of $\alpha$ (the fine-structure constant) and the value of $G$ — two constants that standard physics treats as entirely independent of one another.

\subsection*{Biology as coherence optimisation}

If gravitation is the macroscopic manifestation of coherence leakage at the nuclear scale, one may ask what the PDL says about biological structures, which seem, by contrast, actively to \emph{resist} leakage.

\medskip

The living cell is, in this framework, a remarkable example of hierarchical closure. It maintains an electrochemical gradient across its membrane, selectively exchanges matter and energy with its environment, reproduces its internal organisation with extraordinary fidelity, and actively repairs its own incoherences. Contemporary molecular biology ceaselessly reveals the precision of these mechanisms: chaperone proteins, which help other proteins to find their correct conformation, are literally devices for repairing structural coherence. The mechanisms of DNA replication and correction, which reduce the copying error rate to less than one base in a billion, are biological implementations of the closure-optimisation principle.

\medskip

In this framework, life is not a mysterious phenomenon that escapes the laws of physics: it is the most elaborate form taken by coherence optimisation when closures acquire the capacity to reproduce themselves. Darwinian evolution, far from being a blind and arbitrary process, can be read as a search algorithm in the space of admissible closures: structures that maintain their coherence long enough to reproduce persist; the others dissolve. \emph{Natural selection} is, in this vocabulary, the filter that eliminates closures whose coherence budget is insufficient to ensure their own replication in a given environment.

\medskip

This Darwinian rereading does not weaken the theory of evolution — it anchors it in a more fundamental principle. Genetic variation, drift and selection continue to operate exactly as Darwin described them; what the PDL adds is an explanation of \emph{why} selection is possible, why some structures are more stable than others, and why biological complexity tends to increase: because a higher-level closure can manage its own coherence leakage more economically than a collection of independent closures.

\subsection*{Transcendence: a precise word}

The word \emph{transcendence} often evokes mystery or religion. The PDL proposes a rigorous and sober version of it.

\medskip

A structure is \emph{transcendent} in the PDL sense not because it belongs to another world, but because it cannot fully contain itself. Its coherence overflows. A part of what it is depends on conditions that lie beyond its own boundaries, in the network that surrounds it, in the structures with which it coexists, in a field of global coherence that no individual closure controls.

\medskip

In this sense, transcendence is not the exception; it is the rule. Everything that is sufficiently complex to be interesting is transcendent. Richness and openness go together. One cannot have both complexity and perfect closure at the same time.

\subsection*{Being open without dissolving}

This tension between closure and openness is universal in the PDL. Each structure must find a balance: closed enough to remain itself, open enough to interact, adapt, evolve. Too closed: it isolates itself and can no longer participate in the global dynamics. Too open: it loses its internal coherence and dissolves.

\medskip

We know this tension well in our daily lives. An individual who is too rigid ends up cutting themselves off from the world. An individual with no internal boundaries lets themselves be invaded and no longer knows who they are. Between the two there is a delicate space, that of living coherence, of the structure that maintains itself by adapting. The PDL says that this space is not a psychological metaphor: it is the structural condition of all complex existence, from the proton to the human being.

\subsection*{A universe that cannot close upon itself}

This is not a system that tends towards perfect equilibrium, towards total closure. It is a network of partially open structures that organise themselves with respect to one another, each exporting a fraction of its coherence into a global field that constrains them all.

\medskip

We are all, electrons, protons, human beings, nodes in this fabric. We export coherence. We receive it. We can neither do without it nor dissolve into it. It is this shared condition, this fundamental impossibility of total closure, that the next chapter will help us to reread in terms of causality and time.

% ============================================================
\chapter{Causality Revisited}

Why do things happen in a certain order? Why does a cause precede its effect? Why does time seem to go in only one direction? These questions have been at the heart of physics since Newton. And yet they remain strangely without a deep answer in current theories.

\medskip

Relativity says that time is relative to the observer. Quantum mechanics says that the future is not determined before measurement. But neither really explains \emph{why} there is a direction of time, why causes precede effects, why the world is not simply a frozen block of correlations without order. The PDL offers a different answer, and it begins by questioning the very definition of causality.

\subsection*{Causality as a chain: a convenient but limited image}

The classical image of causality is that of a chain: A causes B, B causes C, C causes D. Each link pushes the next. This image is intuitive, it works well to describe everyday phenomena, and it lies at the basis of all classical physics.

\medskip

But it becomes problematic as soon as one leaves simple situations. In quantum mechanics, an event does not have a single, localised cause: it results from a superposition of states. In general relativity, the temporal order between two events depends on the observer. And at the origins of the universe, the very notions of "before" and "after" lose their meaning.

\subsection*{Causality as global coherence}

In the PDL, causality is not a relation between two isolated events. It is a \emph{global} property: one event is causally linked to another not because it "pushes" it, but because both belong to the same coherence network, which cannot admit them simultaneously in just any order.

\medskip

Imagine a very full agenda: certain appointments cannot be moved without shifting others, not because they touch directly, but because they share common constraints~--- a room, a person, a deadline. It is not the first appointment that "forces" the second; it is the global structure of the agenda that imposes an admissible order. Remove the common constraint, and the order disappears with it. This is exactly what the PDL says: causality is not an arrow between two points, it is a global coherence constraint that is read off the structure of the entire network.

\medskip

What we call a \emph{cause} is a state that opens up a set of coherent transitions. What we call an \emph{effect} is the state that results from the transition most economical in terms of coherence, the one that best preserves the global integrity of the network. The causal order is therefore not an arrow imposed from outside; it is the signature of the way in which a network of closures manages its own maintenance.

\subsection*{ER = EPR: entanglement as shared coherence}

This vision of causality finds an unexpected resonance in one of the most audacious ideas in contemporary physics. In 2013, Juan Maldacena and Leonard Susskind proposed the conjecture $\text{ER} = \text{EPR}$: two quantum-entangled particles would be connected by an Einstein-Rosen bridge~--- that is, by what general relativity calls a wormhole. What seemed to be merely a statistical correlation between two distant particles would in reality be a deep geometrical connection between two regions of space-time.\footnote{J.\,Maldacena \& L.\,Susskind, \textit{Cool horizons for entangled black holes}, Fortschr.\ Phys.\ \textbf{61}, 781 (2013), arXiv:1306.0533. A concrete and calculable realisation of this connection, explicitly deriving the bridge from entanglement in a conformal field theory, was published in November 2024: arXiv:2411.18485.}

\medskip

A word first on what entanglement is, since this is the first time the concept appears in this text. Quantum entanglement is one of the most unsettling phenomena in physics: two particles, once brought into contact, can remain correlated instantaneously regardless of the distance between them. Measuring one seems to "inform" the other, without any signal travelling between them. Einstein called this "spooky action at a distance" and refused to believe it; experiment has proved him wrong.

\medskip

The PDL arrives at the same intuition by an entirely different route. When two coherent closures interact, they do so through \emph{mixed triangles}: relations that span their respective boundaries and create a shared coherence between them. This coupling is not a force acting at a distance; it is a shared structural constraint, exactly as two appointments in the same agenda share a constraint without "touching" one another. Quantum entanglement, in this framework, is not a mysterious phenomenon: it is the relational signature of a coherence shared between boundaries.

\medskip

The convergence between $\text{ER} = \text{EPR}$ and the PDL is not a formal derivation: neither programme follows from the other. It is an independent convergence towards the same central intuition~--- that geometry and information are two readings of the same substrate of coherence~--- reached by entirely different routes. If it withstands testing, it suggests that causality, entanglement, and the geometry of space-time are not three distinct subjects in physics, but three facets of a single principle of relational coherence.

\subsection*{Time as a reading of coherence}

Time, in this framework, is not a backdrop on which events unfold. It is an emergent reading of the order in which coherent configurations succeed one another. When a closure passes from one state to another while maintaining its coherence, it "moves forward in time", not because it follows a universal calendar, but because it realises a sequence of admissible states in the only order possible given its internal constraints.

\medskip

This is not a negation of time; it is an explanation of its origin. Time is the way in which a universe of partial closures tells its own story of coherence maintenance. Where coherence can be preserved, time moves forward. Where it collapses~--- in singularities, at the origins~--- time no longer makes sense, precisely because there is no longer any admissible sequence.

\subsection*{Constants as causal signatures}

The fine-structure constant $\alpha$ does not just say "this is the strength of electromagnetism"; it says: "this is the fraction of coherence that a proton can share with an electron without either losing its integrity". It is a causal constraint: if $\alpha$ were significantly different, atoms as we know them could not maintain their internal coherence, and chemistry, and life, would be impossible.

\medskip

The constants of nature are not arbitrary settings of one universe among other possibles; they are the necessary conditions for a universe of partial closures to remain globally coherent.

\subsection*{Determinism, chance, and freedom}

For simple structures~--- the electron, the proton~--- the sequence of admissible states is highly constrained. The dynamics looks like strict determinism. But for complex structures, those whose relational sea is vast, the number of admissible transitions explodes. There is no longer a single path; there are many, all coherent, all equally possible.

\medskip

What we call chance is the signature of this multiplicity. What we call freedom is the capacity, specific to the most complex structures, to \emph{represent} their own space of admissible transitions and to explore this space actively. Freedom and determinism are therefore not opposed in the PDL; they are two regimes of the same continuum, whose position depends on the internal richness of the structure considered.

\subsection*{Causality without a clockmaker}

What makes this picture particularly striking is the sobriety of the means it employs. The entire edifice of the PDL~--- constants, gravitation, time, causality~--- rests on only two ingredients: integers, which define the discrete structure of the relational network, and a geometric relation, the golden ratio $\phi$, which governs the redistribution of coherence between the core and the active surface of the proton. No adjustable continuous parameters. No constants introduced by hand. Two constraints, and an entire universe follows.

\medskip

The image that emerges from the PDL is that of a universe that does not need a clockmaker, an external law, a prime mover. Causality is not a rule imposed from outside; it is the inevitable consequence of the fact that coherent structures seek to maintain themselves. The order of the world is not programmed; it is \emph{negotiated}, at every moment, by millions of closures that simultaneously manage their own coherence and their inevitable leakage towards what lies beyond them.

\medskip

We are now ready to ask the question that this book has only been preparing from the beginning: what does all this say about \emph{us}? What does it mean to be a human being in a universe governed by this logic of coherence?

% ============================================================
\chapter{We, Vulnerable Instruments}

This chapter is the one that has cost me the most to write. Not for technical reasons, but because it is a matter of saying something honest about us, about what this logic implies if it is true. I wrote it, erased it, rewrote it. What you read here is the version in which I stopped censoring myself.

\medskip

We have followed the logic of coherence from nothingness to the proton, from the proton to the constants of nature, from the constants to causality and time. At each stage, the same law: to exist is to form a stable cycle, to manage a coherence budget, to accept irreducible leakage towards what exceeds us.

\medskip

It is time to ask the question this journey has prepared: what does this logic say about \emph{us}? Not about the electron, not about the proton, but about the human being, with their consciousness, emotions, choices, fragility.

\subsection*{We are high-level closures}

What follows is not a metaphor. It is the direct application of the same irreducible-leakage principle, mathematically demonstrated for the proton, quantified by $\varepsilon_G$, amplified across eighteen levels, to a composite structure of incomparably greater complexity: the human being.

\medskip

A human being is, from the PDL’s point of view, a composite structure of vertiginous complexity. Billions upon billions of electrons and protons, organised into atoms, molecules, cells, organs, systems. Each hierarchical level forms its own closures, manages its own coherence budgets, and exports its own leakage to higher levels and to the environment.

\medskip

But a human being is not only a stack of physical closures. They are also — and this is where the PDL touches something new — a closure that has developed the capacity to represent itself. A structure that can model its own functioning, anticipate its own future states, evaluate its own admissible transitions, and choose among them.

\medskip

It is this leap, from physical closure to reflexive closure, that distinguishes the human being from the electron or the proton. Not a break in the logic of coherence, but its most elaborate expression that we know: a structure that manages its coherence not only by reacting, but by \emph{anticipating}, by \emph{planning}, by \emph{meaning}.

\subsection*{Identity as closure maintained over time}

What is identity? In everyday life, it is what makes you \emph{you} today, tomorrow, in ten years, despite changes of body, ideas, relationships. In the language of the PDL, identity is the maintenance of a coherence pattern over time. Not the rigidity of a structure that never changes — the electron is rigid, but it has no identity in the human sense of the term. Rather, the persistence of an organisation that is reinstantiated at each cycle, slightly modified but recognisable, like a melody that can be played in different keys without ceasing to be itself.

\medskip

Maintaining one’s identity requires continuous effort. It is work of coherence: integrating new experiences without dissolving, adapting without betraying oneself, remaining open without losing the thread. This effort has a cost, a coherence cost in the literal sense of the PDL. And when this cost becomes too high, when contradictions accumulate, when losses exceed the capacity for integration, the closure wavers. What we call an identity crisis, a psychological collapse, or simply exhaustion, can be read as a deficit of internal coherence that the structure can no longer manage alone.

\subsection*{Memory, attention, effort}

Several dimensions of human experience find a precise echo in the vocabulary of the PDL. Not as reductions, but as \emph{resonances} that invite us to think differently.

\medskip

Memory can be read as the capacity to preserve structural regularities beyond a single cycle: coherence patterns that have been stabilised and can be reactivated. To remember is to reinstantiate a past closure in a present context. To forget is to let a pattern degrade to the point that it can no longer be retrieved.

\medskip

Attention looks like an allocation of coherence budget: directing one’s attention to something is to concentrate one’s internal coherence resources on a particular subnetwork, to the detriment of others. That is why attention is limited and fatiguable: it consumes something real.

\medskip

Effort, whether physical or mental, corresponds to the resistance one encounters when trying to modify a well-established closure — to change a habit, learn something new, overcome a fear. The more dense and old the closure, the more expensive it is in coherence to reorganise it. It is inertia, but applied to the mind: no longer the resistance of a physical mass, but the resistance of a tightly woven cognitive pattern.

\subsection*{Freedom: what a complex structure can do}

A structure as complex as the human brain has an internal space of coherent configurations so vast that no computation could enumerate it. Among these configurations, some are more stable, others more fragile. Some open up new possibilities, others close doors. Human freedom is the navigation in this space, guided by experience, constrained by biology and history, but real, because the space itself is real.

\medskip

This freedom is not infinite. It is bounded by the available coherence budget, by closures already established, by irreducible leakage towards an environment one does not control. But it is vast enough that what we call a choice, a deliberation, a decision, a commitment, is not an illusion. It is real navigation in a real space of admissible possibilities.

\subsection*{What medicine might read differently}

Contemporary medicine, especially in its most advanced branches, faces a paradox: the more precise it becomes in the molecular description of diseases, the more it struggles to grasp the patient as a coherent whole. The PDL offers here a complementary perspective, without claiming to replace medical biology.

\medskip

A disease, within the PDL, can be read as a disturbance of the coherence budget of a high-level closure. Cancer, for example, is a situation in which individual cells have “unlearned” how to respect the coherence constraints of the tissue to which they belong: they optimise their own local closure at the expense of the organism’s global coherence. It is precisely the situation described in the analogy of the group of friends: each member seeks to maximise their own relations to the detriment of collective coherence.

\medskip

Autoimmune diseases illustrate another regime: the immune closure confuses its own structures with foreign ones, generating internal triangular incoherence that turns against the organism. And neurodegenerative diseases such as Alzheimer’s can be read as a progressive degradation of the coherence pattern that constitutes identity: memory, interpreted as the capacity to reinstantiate past coherence patterns, dissolves as the underlying structures lose their integrity.

\medskip

These rereadings are not diagnoses; they are invitations to think differently about what medicine already does. Cell therapies, immunotherapies and precision medicine all seek, in their own way, to restore or strengthen coherence constraints that allow an organism to maintain itself. The PDL suggests that this medical project and the PDL’s fundamental physical project share a deep structure: in both cases, it is a matter of understanding how complex closures manage their coherence budget, and how to repair it when it is disturbed.

\subsection*{Vulnerability: the structural price of complexity}

The result here is structural, not metaphorical. In the PDL, coherence leakage $\varepsilon$ grows with hierarchical complexity: the richer a structure is, the larger its incompressible share of openness to the outside. The human being, the most complex composite structure we know, is the most extreme consequence of this.

\medskip

A human being is the most elaborate composite structure we know. Their coherence leakage is therefore immense, not in the physical sense of gravitational exponentiation, but in the existential sense: a considerable portion of what maintains a human being in coherence lies \emph{outside} them. In other people who recognise them. In the institutions that organise them. In the ecosystems that nourish them. In the languages that allow them to think. In the shared memories that give them a history.

\medskip

Remove these supports — total isolation, loss of recognition, destruction of the social fabric — and the human closure begins to degrade. This is not a metaphor: it is a well-documented clinical observation. Human vulnerability is not a contingent weakness that one could eliminate with more will or technique. It is the direct and inevitable consequence of our complexity. We are vulnerable because we are rich, because our coherence is so dense and so hierarchical that it cannot be maintained without a vast and diversified external support network.

\subsection*{Relation as a condition of existence}

If human coherence depends structurally on the outside, then relation is not a luxury or an ornament of life: it is a condition of existence\footnote{This statement has a precise technical content in the PDL: a structure whose coherence leakage exceeds its internal absorption capacity can only be maintained by importing coherence from its environment. Relational dependence is not a psychological property — it is a structural constraint on any high-complexity closure.}. To be in relation — with other people, with nature, with ideas, with a tradition — is to feed one’s own closure with external coherence. It is what allows the structure to maintain and develop itself.

\medskip

The PDL offers here, cautiously, a structural reading of what great ethical and spiritual traditions have always intuited: that the human being is fundamentally relational, that their dignity is inseparable from their fragility, and that to care for others is also to care for the coherence fabric that maintains us all in existence.

\subsection*{Instruments, not victims}

The guiding sentence of this paper calls the human being “the most vulnerable of instruments”. An instrument, in the PDL sense, is not passive: it is a structure through which a logic is realised, a form through which something universal finds a particular, local, irreplaceable expression.

\medskip

The electron is an instrument: through it, the logic of minimal coherence is realised for the first time in the physical world. The proton is an instrument: through it, that same logic builds its first complex architecture. And the human being is an instrument — the most elaborate, the richest, the most fragile — through which the logic of coherence reaches the level at which it can \emph{reflect upon itself}.

\medskip

This is, perhaps, what is most singular in the human condition: we are not only structures that maintain themselves, we are structures that \emph{know} that they maintain themselves, that can contemplate the logic that runs through them, that can choose to participate in it consciously. We are the place where the universe begins to understand itself.

\medskip

This does not place us above the logic of coherence; it plunges us more deeply into it. Our reflexivity is not an escape route: it is a responsibility. That of managing our coherence with care — our own, and that of the relational fabric on which we depend and which we help to weave.

\medskip

Let us take a concrete example. Someone loses a loved one. In ordinary language, we say they are “broken”. In the language of the PDL, one could say that their identity closure — that pattern of relations, memories, habits, projects, that constitutes their self — has undergone a massive disturbance. Part of the coherence network that defined them has disappeared with the other. They must rebuild themselves: find a new balance, redistribute their coherence budget, reorganise their internal relations. This is painful, not metaphorically but structurally: it is the real cost of a deep reorganisation of a complex closure. And what makes this reconstruction possible is precisely leakage: the incompressible openness towards others, towards the world, which prevents the closure from shutting itself in on its own suffering.

\medskip

Our vulnerability is not a flaw. It is the condition of our openness. And our openness is the condition of our existence.

% ============================================================
\chapter[Conclusion --- Learning to Read the Logos]{Conclusion \\ Learning to Read the Logos}

We set out from nothingness. Not as a poetic backdrop, but as a logical challenge: what can exist without presupposing anything else? The PDL's answer is radically sober: a difference that can repeat itself, a cycle that returns to itself, a pulse. That is all. And yet, from this single principle, an entire architecture has emerged: the electron, the proton, the constants of nature, gravitation, time, and even the human condition.

\subsection*{A programme that knows how to dialogue}

The PDL does not advance in isolation. It has engaged in a formal dialogue with another fundamental research programme: the Theory of Objectivity (TO), which seeks to identify the minimal logical conditions that any admissible universe must satisfy in order to allow the persistent existence of distinct objects for several observers.\footnote{\textit{A Modal Dialogue Between the Projective Dynamic Logo and the Theory of Objectivity}, C.~Laubscher, Zenodo, 2026.}

\medskip

This dialogue is not an academic curiosity. It has produced concrete refinements within the PDL programme: a more rigorous definition of active surfaces, a more explicit treatment of boundaries between a structure and its environment, and deeper reflection on what it means to "exist for more than one observer". These refinements are not cosmetic; they have made certain claims of the programme more robust and more sharply delimited.

\medskip

What is worth retaining is that a research programme's capacity to dialogue with its intellectual neighbours without dissolving~--- without losing its internal coherence~--- is itself a form of test. A programme that is too fragile collapses at the first critical glance. A programme that is too rigid closes itself off and no longer progresses. In this exchange, the PDL has shown that it can do what its own structures do: remain open without dissolving. Accept leakage towards what exceeds it, and emerge from it more precise.

\subsection*{What the PDL is not}

The PDL is not a finished theory that would replace current physics overnight. It is not a philosophical doctrine that claims to explain everything. It is a research programme: a set of precise mathematical constructions, testable hypotheses, conjectures honestly identified as such, and clearly posed open questions. This epistemic honesty is, in our view, one of its most precious qualities.

\subsection*{What the PDL changes}

It changes the way we think about existence: to exist is no longer a mysterious property, it is a performance, an active maintenance of coherence under constraint. It changes the way we think about the constants of nature: these numbers are no longer arbitrary; they encode the conditions under which a universe of partial closures can remain globally coherent. It changes the way we think about vulnerability: to be fragile is no longer a defect to be corrected, it is the structural condition of all complex existence. It changes, perhaps, the way we think about transcendence: no longer a separate beyond, but the structural impossibility of fully containing oneself.

\medskip

It also changes the way we think about what counts as a serious theory. A serious theory is not only a theory that predicts correct things; it is a theory that says precisely what would destroy it.

\subsection*{The PDL knows what could refute it. And it says so.}

First scenario: the constants $\alpha$ and $G$ are locked, in the PDL, by a precise proton architecture defined by five integers~--- $n_u = 24$, $n_d = 28$, $r_\text{val} = 930$, $R_\text{mer} = 10\,087$, $R_\text{tot} = 11\,017$. This architecture is unique: changing one of these integers, even by a single unit, destroys agreement with the measured values. If future measurements of $\alpha$ or $G$ were to deviate from the corridor defined by this architecture~--- not by a per cent, but by a minute fraction, for the window is extraordinarily narrow~--- the bridge between the two constants would collapse. This would not be a correctable detail. It would be a refutation.

\medskip

Second scenario: the PDL asserts that the $(24, 28, 930, 10\,087, 11\,017)$ architecture is the unique solution that simultaneously satisfies the programme's constraints. If someone were to find an alternative architecture~--- different integers, a different logical proton, a different hierarchical closure~--- that reproduces the same constants with the same precision, then the claim to uniqueness would fall. The programme would remain interesting, but it would lose one of its strongest assertions.

\medskip

Two more recent results of the programme deserve mention, because they illustrate two very different ways in which the framework can be tested.

\medskip

The first concerns light and cavities. In ordinary physics, a closed box~--- a cavity~--- can contain light. If one counts how many different vibration modes this light can adopt as a function of frequency, one obtains a well-known law: the number of modes grows like the square of the frequency, written $\rho(\nu) \propto \nu^2$. This law underlies blackbody radiation and, ultimately, all quantum field physics. The PDL has shown that a cavity built \emph{solely} from finite signed graphs, without presupposing the prior existence of space or of a wave equation, spontaneously supports periodic modes whose density follows exactly this same law $\rho(\nu) \propto \nu^2$, in three dimensions\footnote{\textit{Discrete Cavity Modes and Logical Density of States in the PDL Framework}, C.~Laubscher, Zenodo, 2026.}. This result was not engineered to reproduce this law; it is its consequence. It is an independent, non-circular validation: the same framework that selects the electron and the proton reproduces, without additional parameter, a law that physics has regarded since Planck as a fundamental property of three-dimensional space. This suggests that the three-dimensional structure of space itself could emerge, in this programme, as a consequence of relational coherence constraints, rather than as a presupposed backdrop.

\medskip

The second result concerns one of the oldest challenges in theoretical physics: the unification of general relativity and quantum mechanics. These two theories describe the world with extraordinary precision, each in its own domain, but their fundamental languages are incompatible. The PDL sketches a framework in which this incompatibility could find a natural resolution: the $(4,6)$ structures that model the electron can be read as germs of Dirac spinors, while the emergent metric arising from coherence cycles provides an analogue of Einsteinian curvature\footnote{\textit{Towards an Einstein--Dirac Unification from the PDL Framework}, C.~Laubscher, Zenodo, 2026.}. Preliminary consistency checks~--- hydrogenic limit, gravitational redshift~--- have been carried out with encouraging results. It is essential to be clear about the status of this result: it is a structural sketch, not an accomplished unification. The PDL does not claim to have completed unification. But it is the only programme to date to have derived~--- from a single principle of logical closure, without a free parameter~--- both the form of the electron, the architecture of the proton, the great constants of nature, and a structural common language for general relativity and quantum mechanics. This is not a correction. This is not a reformulation. It is a change of foundations: and if the foundations hold, what they carry holds with them.

\medskip

These two scenarios are not causes for fear; they are guardrails. They show that the PDL is not built to withstand everything. It is built either to be correct, or to be shown wrong in a precise, irrefragable way. That is the very definition of a serious scientific programme.

\subsection*{Energy, mechanical production, and the cost of closure}

The PDL has concrete implications for how we think about the production and transformation of energy. In the classical framework, energy is a conserved quantity that transforms from one form into another; thermodynamics describes its limits. In the PDL framework, energy is a coherence budget: the amount of internal organisation that a structure can maintain and transfer.

\medskip

This rereading has consequences for our understanding of machines. A heat engine, seen through the PDL, is a closure that exploits a coherence gradient between two reservoirs (hot and cold) to produce ordered work. Carnot efficiency~--- the maximal efficiency of a heat engine~--- is, in this vocabulary, the constraint on the fraction of coherence that a closure can transfer from a disorganised reservoir (hot) to a more organised closure (cold) without violating the second law of thermodynamics. What Boltzmann described as entropy~--- and what the PDL rereads as a measure of coherence leakage distributed over a very large number of degrees of freedom~--- is precisely the macroscopic translation of the PDL leakage parameter $\varepsilon$, aggregated over billions upon billions of elementary closures.

\medskip

Recent developments in energy~--- nuclear fusion, hydrogen fuel cells, superconducting materials~--- can all be read as attempts to build closures capable of maintaining or transferring coherence with minimal leakage. The challenge of thermonuclear fusion is precisely to maintain a plasma at more than one hundred million degrees in a sufficiently coherent state for fusion reactions to occur: it is a closure problem under the most extreme conditions. The recent advances of the ITER reactor and the promising results of inertial-confinement fusion (notably the net-positive energy records obtained at NIF in 2022--2023) illustrate concretely how fundamental the problem of plasma coherence management is, and how difficult it is to maintain an open closure without dissolving it.

\subsection*{The PDL in the face of recent discoveries}

A serious research programme does not live in a vacuum: it must confront the most recent experimental discoveries, not to \emph{explain} them prematurely, but to test whether its conceptual framework is compatible with them.

\begin{itemize}

\item \textbf{The images of the black holes M87* and Sgr A*.} In 2019 and 2022, the Event Horizon Telescope network produced the first direct images of an event horizon~--- that of the supermassive black hole in galaxy M87, then that at the centre of our own galaxy. These images show a luminous ring surrounding a shadow~--- exactly what Einstein's general relativity predicts. In the PDL framework, a black hole is a region where the density of closures is so extreme that local coherence leakage can no longer propagate outwards: the event horizon\footnote{Boundary beyond which no signal, not even light, can escape from a black hole.} is the boundary beyond which no coherence can be exported. The preliminary sketch of the Einstein--Dirac interface in the PDL~--- explicitly identified as speculative~--- thus finds a first geometric confrontation with observation.

\item \textbf{Gravitational waves and LIGO.} Since the first direct detection of gravitational waves by LIGO\footnote{Laser Interferometer Gravitational-Wave Observatory; an American detector that made the first direct observation of gravitational waves in 2015.} in 2015, more than a hundred events have been recorded~--- mergers of black holes and neutron-star pairs. These waves are, in classical physics, ripples in space-time curvature. In the PDL, they would be collective propagations of coherence leakage: perturbations of the global coherence field generated by extremely violent events involving very massive structures. The reinterpretation of $G$ as the result of hierarchical filtering across eighteen levels is precisely the kind of prediction that precision measurements of gravitational waves could, in the long run, test.

\item \textbf{Proton-mass measurements and the constants crisis.} Recent ultra-precise measurements of the proton mass (notably via Penning traps) and ongoing discussions about slight tensions between different measurements of the Hubble constant (the ``Hubble tension''\footnote{A statistical disagreement between the two main methods of measuring the universe's expansion rate; one based on the cosmic microwave background, the other on Cepheids and supernovae.}) raise questions about the stability of fundamental constants. If the PDL is right in tying $\alpha$ and $G$ to a precise proton architecture, then any temporal or spatial variation of these constants would be the signature of a variation in the architecture itself~--- which is difficult to conceive within the PDL, but nonetheless constitutes a valuable coherence test.

\item \textbf{Synthetic biology and the origins of life.} Advances in synthetic biology~--- in particular the laboratory creation of minimal cells capable of division (JCVI experiments since 2010, culminating in \textit{JCVI-syn3.0} in 2016 and \textit{syn3A} in 2021)~--- raise the question of what is minimally required for a biological structure to be ``alive''. In the PDL framework, a minimal viable cell is the smallest biological closure capable of maintaining its internal coherence \emph{and} reproducing itself. The discovery that \textit{syn3.0} contains about 473 genes, of which the function of a third remains unknown, suggests that the space of minimal biological closures is not yet fully charted~--- and that the PDL could offer structural criteria for understanding why some configurations are viable and others are not.

\item \textbf{Non-equilibrium physics.} The work of Jeremy England on adaptive dissipation shows that physical systems subject to an external energy flux tend spontaneously towards configurations that maximise their dissipation~--- what he sometimes calls a ``thermodynamics of evolution''. This idea is structurally close, though formulated in a different language, to the PDL principle that the most stable closures are those that best manage their coherence leakage to the environment. A formal dialogue between these two approaches remains to be built, but the conceptual convergence is striking.

\end{itemize}

\subsection*{Returning to the sentence}

Let us return, one last time, to the sentence that opened this text:

\begin{quote}
\itshape Whoever we may be, from the atoms that form our molecules to our most intimate thoughts, we are nothing but the necessary expression of a simple and inexorable logic of coherence. A logic of causality that fulfils itself at every pulse. It is the Origin, the engine and the essence of the universe. The human being is merely its most vulnerable instrument.
\end{quote}

This sentence is not a slogan. It is a working hypothesis, formulated precisely, anchored in rigorous mathematical constructions, open to refutation. But if it is correct, even partially, then it says something profound about what it means to exist. It means that every time you maintain a coherent thought, every time you keep a promise, every time you repair a damaged relationship, you participate in the same project as the electron that loops upon itself. The vocabulary changes, the scale changes, the richness changes, but the deep structure is the same: maintaining coherence, accepting leakage, remaining open without dissolving.

\subsection*{An invitation, not a doctrine}

The PDL is a young, ambitious, incomplete programme. It needs to be tested, criticised, refined~--- by physicists, philosophers, mathematicians, by all those who take an interest in it, who bring questions its authors have not yet thought to ask.

\medskip

If you are a physicist or mathematician, the technical publications, freely available on \href{https://zenodo.org}{Zenodo}, await you with their proofs, conjectures, and open questions. If you are a philosopher, you will find in them a precise and contestable ontology. And if you are simply someone who wonders what it means to live in this universe, then we hope this paper has given you a few new tools for posing that question more clearly.

\medskip

For perhaps that is what, in the end, it means to learn to read the Logos: not to decipher a sacred text, not to master a definitive theory, but to learn to recognise, in everything that endures, the silent work of coherence, and in every fragility, the signature of an irreducible opening towards what exceeds us.

\medskip

I do not know whether the PDL is correct. I know that it is coherent, that it is testable, and that it has not yet been found wanting. That is all one can ask of a research programme at this stage. What I hope is that it will at least have made you want to pose the question differently.

\chapter{If This Vision Is Correct\texorpdfstring{\,\ldots}{...}}

The PDL is a research programme, not a doctrine. But every serious programme must answer a precise question: \emph{what does it predict that we do not yet know?} This chapter does not recapitulate what has been established. It looks ahead: towards the tests that await completion, the dialogues that remain to be built, and the domains where the logic of coherence, if it proves fundamental, compels us to reformulate questions we believed settled.

\subsection*{What the PDL Predicts --- and Has Not Yet Been Measured}

The strength of a scientific programme is measured not only by what it explains from the past, but by what it commits to predicting for the future. The PDL formulates several precise predictions, some of which are today technically within reach.

\medskip

\textbf{The fine-structure constant $\alpha$ and the proton-to-electron mass ratio $\mu$.} The PDL locks both $\alpha$ and the gravitational constant $G$ simultaneously from a unique architecture defined by five integers $(24, 28, 930, 10\,087, 11\,017)$. This entails an implicit prediction about the ratio $\mu = m_p / m_e$\footnote{The ratio $\mu = m_p/m_e \approx 1836$ is the number of times the electron ``fits'' into the mass of the proton. Its precise value has been a puzzle since the beginnings of atomic physics: nothing in the standard model explains why this number is exactly what it is.}: if the architecture is correct, this ratio cannot be even slightly different from what we measure, for it is constrained by the same logical closure. Measurements of $\mu$ via precision molecular spectroscopy\footnote{Precision spectroscopy consists of measuring very accurately the frequencies of light absorbed or emitted by molecules or atoms. By measuring these frequencies to a precision of around $10^{-12}$, one can test whether the fundamental constants take exactly their expected values.}~--- notably on H$_2^+$ and HD$^+$ in ion traps\footnote{An ion trap is a device that holds charged atoms or molecules suspended in a vacuum using electric and magnetic fields, allowing them to be studied without disturbance from collisions.}~--- currently reach precisions of the order of $10^{-11}$ to $10^{-12}$. Any deviation from the value predicted by the architecture would constitute a direct refutation.

\medskip

\textbf{The topological signature of the electron in colliders.} The PDL models the electron as a minimal stationary closure of type $(4,6)$\footnote{The notation $(4,6)$ denotes a signed graph with 4 positive vertices and 6 edges: the smallest structure capable of sustaining a stable coherence cycle according to the PDL's rules. It is not a geometric object in ordinary space, but an abstract relational structure.}~--- the smallest signed graph capable of sustaining a stable coherence cycle. This structure is not a mere point: it possesses an internal topology that could, in principle, leave traces in high-energy scattering processes. At the LHC\footnote{The Large Hadron Collider (LHC) at CERN, near Geneva, is the most powerful machine ever built for studying matter at the subatomic scale: it accelerates protons to speeds approaching that of light and collides them, revealing their constituents.}, the ATLAS and CMS experiments measure angular and polarisation distributions in processes involving electrons at energies of the order of a TeV\footnote{A TeV (tera-electronvolt) is a unit of energy in particle physics: it is roughly one trillion times the energy of a visible-light photon. At this scale, the internal structure of particles, if any exists, becomes in principle detectable.}. If the $(4,6)$ structure possesses a geometric signature in these processes~--- for instance a slight deviation in angular distribution functions or in polarimetric Stokes parameters\footnote{The Stokes parameters are quantities that describe the polarisation state of light or particles: they allow one to measure whether particles produced in a collision preferentially ``spin'' in one direction.}~--- it would in principle be detectable. The PDL has not yet derived the precise differential cross-section\footnote{The differential cross-section is a measure of the probability that a particle is scattered in a given direction during a collision. It is the quantity that particle-physics experiments measure most directly.} that this structure implies: this is an open problem explicitly identified within the programme\footnote{\textit{PDL Global Mapping of Structures, Results and Open Problems}, v7, C.~Laubscher, 2026, available on \href{https://github.com/laubscher-lab/PDL-framework}{GitHub}.}.

\medskip

\textbf{The proton as a structural analogue of black holes.} One of the most unexpected predictions of the PDL concerns the parallel between the structure of the proton and that of a black hole. In both cases, the PDL describes a closure whose active surface encodes the total information contained in the volume: for the proton, this is the Mersenne surface\footnote{The Mersenne surface, in the PDL, is the layer of cycles forming the active boundary of the proton: the level at which the proton's internal information ``dialogues'' with the outside. Its name derives from the fact that the number $R_\text{mer} = 10\,087$ is linked to a combinatorial structure of Mersenne type.} defined by $R_\text{mer} = 10\,087$ cycles; for a black hole, it is the event horizon\footnote{The event horizon of a black hole is the imaginary boundary beyond which nothing~--- not even light~--- can escape. It is not a solid surface: it is a point of no return defined by the geometry of space-time.}. If this parallel is structurally founded, it implies a precise prediction: the ratio between the entropy of the proton\footnote{The Bekenstein--Hawking entropy is a measure of the number of internal states of a black hole, proportional to the area of its horizon. The PDL's conjecture is that the proton obeys an analogous scaling law, which would constitute a deep link between particle physics and gravitation.} (in the generalised Bekenstein--Hawking sense) and its charge radius should follow the same scaling law as for black holes. Ultra-precise measurements of the proton's charge radius~--- notably via laser spectroscopy of muonic hydrogen\footnote{Muonic hydrogen is an artificial atom in which the ordinary electron is replaced by a muon~--- a particle 200 times heavier. Because the muon orbits far closer to the proton, it is much more sensitive to its exact size. In 2010, measurements on this atom revealed a proton radius slightly different from values obtained with ordinary hydrogen: this is what became known as the ``proton radius puzzle''.}, which revealed in 2010 the famous ``proton radius puzzle''~--- could constitute an indirect test of this parallel. If the PDL predicts a precise value of the charge radius from its architecture, that prediction is computable and comparable to current measurements.

\subsection*{Cosmology: the Hubble Constant, Dark Matter, Dark Energy}

\textbf{The Hubble tension reread as an architectural signal.} The ``Hubble tension'' refers to the persistent disagreement between two independent measurements of the rate of expansion of the universe\footnote{The universe is expanding: galaxies are moving away from one another. The Hubble constant $H_0$ measures the speed of this expansion: it indicates how many km/s a galaxy recedes for each megaparsec of distance (one megaparsec is approximately 3.26 million light-years).}: the cosmic microwave background method\footnote{The cosmic microwave background (CMB) is the fossil radiation emitted 380,000 years after the Big Bang, when the universe had cooled sufficiently to become transparent. Its temperature and fluctuations allow the cosmological parameters to be reconstructed with very high precision.} gives $H_0 \approx 67.4$ km/s/Mpc, while the local distance ladder (Cepheids\footnote{Cepheids are stars whose luminosity oscillates periodically. Because this period is linked to their intrinsic luminosity, they serve as ``standard candles'' for measuring distances in the universe.}, Type Ia supernovae) gives $H_0 \approx 73$ km/s/Mpc. This disagreement, exceeding $5\sigma$\footnote{In scientific statistics, a $5\sigma$ discrepancy means the probability that the disagreement is due to chance is less than one in 3.5 million. This is the conventional threshold for claiming a ``discovery'' in particle physics.}, has resisted all known systematic corrections for more than a decade. In the PDL, $G$ is the result of hierarchical filtering across eighteen levels of coherence-leakage amplification. If this filtering is not perfectly uniform at cosmological scales~--- if there exists a slight dependence of $G$ on the local density of closures~--- then $H_0$ could indeed take slightly different values depending on the density regime being probed. This is not an explanation of the tension, but a structural prediction: if the PDL is correct, the Hubble tension should \emph{persist} and its amplitude should be correlated with the scale at which $G$ is measured.

\medskip

\textbf{Dark matter as a population of non-luminous closures.} The PDL does not predict a priori the existence of weakly interacting massive particles (WIMPs\footnote{WIMPs (\textit{Weakly Interacting Massive Particles}) are hypothetical dark-matter particles: as massive as atomic nuclei, but interacting so weakly with ordinary matter as to be nearly undetectable. They are one of the most extensively studied candidates for explaining dark matter.}). But it does predict the existence of a spectrum of stable closures: structures that satisfy the programme's coherence constraints without necessarily coupling to photons. A stable non-luminous closure would, in this vocabulary, be a structure whose active surface supports no electromagnetic modes~--- that is, a closure whose architecture differs from that of the electron and the proton, but is equally constrained. The PDL therefore generates a qualitative prediction: if dark matter exists, it should possess precise topological properties~--- notably a well-defined surface entropy and a mass linked to an integer number of coherence cycles. This prediction remains qualitative for now; making it quantitative is an open problem.

\medskip

\textbf{Dark energy as the energy of the relational vacuum.} In standard cosmology, dark energy is modelled by a cosmological constant $\Lambda$\footnote{The cosmological constant $\Lambda$ was introduced by Einstein in 1917 to maintain a static universe~--- he later withdrew it, calling it his ``greatest blunder''. It was reintroduced in 1998 when the accelerating expansion of the universe was discovered. Its measured value is $10^{120}$ times smaller than what quantum field theory predicts: this is one of the most glaring unresolved problems in modern physics.}~--- a vacuum-energy term whose measured value is $10^{120}$ times smaller than what particle physics predicts: the ``ultraviolet catastrophe'' of modern cosmology. In the PDL, the vacuum is not empty: it is populated by minimal closures in dynamic equilibrium, whose aggregated coherence leakage constitutes a residual energy of the global relational network. If this residual energy is proportional to the total number of elementary closures in the observable universe~--- a finite and computable number within the PDL framework~--- then the value of $\Lambda$ could be derived rather than postulated. This calculation has not yet been carried out: it constitutes one of the most ambitious open problems of the programme.

\medskip

\textbf{The Dark Energy Survey and large-scale structure.} In January 2026, the DES collaboration published the combined analysis of its full six years of data~--- the most precise mapping to date of the distribution of matter and dark energy at large scales\footnote{DES Collaboration, \textit{Dark Energy Survey scientists release analysis of all six years of data}, January 2026.}. The results show a slight tension with the standard $\Lambda$CDM model\footnote{The $\Lambda$CDM (\textit{Lambda Cold Dark Matter}) model is the standard cosmological model: it describes a universe composed of approximately 5\,\% ordinary matter, 27\,\% cold dark matter, and 68\,\% dark energy represented by $\Lambda$. It is the most precise model available, but leaves unanswered the fundamental questions about the nature of its two dominant components.} in the gravitational shear parameters\footnote{Gravitational shear (\textit{weak gravitational lensing}) is the slight distortion of images of distant galaxies caused by matter~--- ordinary and dark~--- lying along the line of sight. By statistically measuring these distortions across millions of galaxies, one can map the distribution of mass in the universe.}, suggesting that either dark energy is not a pure constant, or gravitation behaves slightly differently from Einstein's predictions at very large scales. Within the PDL framework, these tensions are precisely what one would expect if $G$ results from hierarchical filtering across eighteen levels: at scales where the density of elementary closures becomes sufficiently low, the filtering might no longer be uniform, producing an \emph{effective gravitation} that varies slightly with the density regime. The PDL therefore predicts that these DES tensions should \emph{persist and amplify} with future surveys~--- notably DESI\footnote{DESI (\textit{Dark Energy Spectroscopic Instrument}) is a spectrograph installed on the Mayall Telescope in Arizona, capable of measuring the distance of 5,000 galaxies simultaneously. Euclid is an ESA satellite launched in 2023 to map the geometry of the dark universe. The Rubin Observatory in Chile will begin in 2025 the largest photographic survey ever undertaken.}, Euclid, and the Rubin Observatory~--- and that they should be correlated with similar anomalies in the measurement of $H_0$. This is not an explanation: it is a structural prediction.

\medskip

\textbf{Direct dark-matter detection and the neutrino floor.} Direct dark-matter detection experiments~--- notably LZ (\emph{LUX-ZEPLIN})\footnote{LUX-ZEPLIN is a dark-matter detector installed 1.5 km underground in a gold mine in South Dakota. It contains 10 tonnes of ultra-pure liquid xenon: if a dark-matter particle traverses this volume, it should produce a tiny flash of light. The depth shields the detector from cosmic-ray interference.}, XENONnT, and PandaX-4T~--- reached a historic threshold in 2025: the \emph{neutrino floor}\footnote{The neutrino floor is the limit below which dark-matter detectors can no longer distinguish a potential dark-matter signal from the background produced by solar neutrinos~--- near-undetectable particles emitted continuously by the Sun. Reaching this floor means that current detectors have attained their fundamental limit for WIMP searches.}, the regime in which the irreducible background is constituted by solar neutrinos interacting via coherent elastic neutrino-nucleus scattering (CE$\nu$NS). LZ in particular excluded WIMPs in the 3--9 GeV/c$^2$ window with record-setting limits, with no positive signal\footnote{LZ Collaboration, \textit{Dark matter search sets new limits and achieves neutrino floor}, December 2025.}. In the language of the PDL, this absence is structurally significant: if dark matter is a population of stable non-luminous closures, as the programme predicts, then these closures should not interact with ordinary nuclei via the exchange of standard gauge bosons\footnote{Gauge bosons are the particles that ``carry'' the fundamental forces: the photon for electromagnetism, the W and Z bosons for the weak force, gluons for the strong force. A non-luminous closure, in the PDL, would be a structure that couples to none of these carriers.}. The absence of WIMPs in the explored windows is therefore \emph{compatible} with the PDL prediction of a dark matter of topological rather than particulate nature. The PDL's precise prediction~--- that the mass of stable non-luminous closures should be linked to an integer number of coherence cycles, and not freely distributed across a continuous spectrum~--- implies that they could reside at very precise masses, still beyond the reach of current detectors, or conversely well below the GeV\footnote{A GeV (giga-electronvolt) is a unit of mass-energy in particle physics: the proton weighs approximately 0.938 GeV. Current detectors search primarily for WIMPs between 1 and 1,000 GeV. The sub-GeV region (below 1 GeV) remains largely unexplored.}. This is a test that future sub-GeV detection programmes~--- notably DarkSPHERE and the phonon detectors at SNOLAB\footnote{SNOLAB is a Canadian underground laboratory situated 2 km below the surface in a nickel mine in Ontario. Phonon detectors measure vibrations of the crystal lattice produced by a collision: they can in principle detect particles far lighter than liquid-xenon detectors.}~--- could address as early as 2026--2027.

\subsection*{Tests in Nanotechnology and Synthetic Biology}

\textbf{Nano-living structures and minimal closures.} Recent advances in synthetic biology~--- notably the minimal cells JCVI-syn3A\footnote{JCVI-syn3A is a synthetic bacterial cell created by the J. Craig Venter Institute: it contains only 473 genes, compared to several thousand in an ordinary bacterium. It is the simplest artificial cell ever built that is capable of growing and dividing. A third of its genes have as yet unknown functions, suggesting that we do not yet fully understand what is minimally necessary for life.} (473 genes, of which a third have unknown function)~--- raise the question of what is \emph{minimally} necessary for a biological structure to be viable. The PDL predicts that viability is a closure property: a minimal viable cell is the smallest biological closure capable of maintaining its internal coherence \emph{and} reproducing itself. This prediction is testable: if synthetic cells are constructed with decreasing numbers of functional components, the PDL predicts the existence of a \emph{hard threshold} of viability~--- not a progressive degradation, but an abrupt transition between ``coherent closure'' and ``dissolution''. Current genome-reduction experiments do not yet clearly show whether this transition is abrupt or gradual: this is precisely what the PDL predicts, and what synthetic biology could test.

\medskip

\textbf{Artificial nanostructures and the golden ratio.} The PDL has shown that the active surface of the proton spontaneously gives rise to the golden ratio $\varphi = (1+\sqrt{5})/2$\footnote{The golden ratio $\varphi \approx 1.618$ is a mathematical number that appears in widely varying contexts: the spiral of shells, the arrangement of sunflower seeds, the proportions of the Parthenon. In the PDL, its appearance is not a sign of mysterious harmony, but of a combinatorial constraint: it is the ratio that minimises coherence leakage from a surface under integer constraints.} as the ratio between surface cycles and total cycles\footnote{\textit{Emergence of the Golden Ratio in the PDL Proton Surface}, C.~Laubscher, 2026.}. This emergence is not aesthetic: it is structural, linked to the stability conditions of a closure under combinatorial constraint. If the golden ratio is a general property of optimally stable closures~--- and not an accident of the proton~--- then artificial nanostructures designed to maximise their mechanical stability under constraint should converge towards geometric ratios close to $\varphi$. Tests on fullerenes\footnote{Fullerenes are molecules of pure carbon formed from pentagons and hexagons: the best known, buckminsterfullerene C$_{60}$, resembles a football. Their remarkable stability makes them natural candidates for testing whether stable structures under constraint converge towards precise geometric ratios.}, carbon nanotubes, and DNA origami structures\footnote{DNA origami is a technique that folds DNA strands into precise three-dimensional structures: nanocubes, spheres, octahedra. These structures are fully programmable and allow geometric predictions to be tested at the nanometric scale.} could verify whether this property is universal or specific.

\subsection*{Quantum Computing: Sustaining Openness Without Dissolving}

Quantum computing is, to date, the technological domain in which the tension between coherence and leakage is most visible, most costly, and most precisely measured. A quantum computer calculates correctly only if its qubits\footnote{A qubit (\textit{quantum bit}) is the basic unit of quantum information. Unlike a classical bit, which is either 0 or 1, a qubit can exist in a superposition of both states simultaneously~--- until it is measured. This property is both the strength of quantum computing and its fragility: the slightest coupling with the environment destroys the superposition.} remain in coherent superposition throughout the duration of the computation. Decoherence~--- the involuntary coupling of the system with its environment~--- is the principal enemy. This problem is not an engineering detail: it is, in the language of the PDL, the fundamental problem of sustaining an open closure without dissolving it.

\medskip

\textbf{Qubits as partial closures.} In the PDL, a qubit is not a point-like object with two states: it is a partial closure whose active surface can reside in a superposition of two coherence configurations. Decoherence is then precisely the leakage $\varepsilon$ of the PDL~--- the impossibility for a closure to remain fully closed upon itself in the presence of an environment. What quantum engineers call ``coherence time'' $T_2$\footnote{The time $T_2$ (phase coherence time) measures how long a qubit can sustain its superposition before interaction with the environment destroys the phase information. The best current superconducting qubits achieve $T_2$ values of the order of a few milliseconds: a record, but still insufficient for the most ambitious algorithms.} is, in this vocabulary, the duration for which a partial closure can maintain its configuration before leakage prevails. This rereading is not merely metaphorical: it predicts that $T_2$ should depend on the topology of the coupling between the qubit and its environment in a precise manner, linked to the structure of the PDL's active surfaces.

\medskip

\textbf{Quantum error-correcting codes as implementations of minimal leakage.} Quantum error-correcting codes~--- surface codes\footnote{A surface code is an error-correcting code that encodes one logical qubit in a network of several physical qubits arranged on a two-dimensional grid. By measuring combinations of neighbouring qubits (without measuring the individual qubits), errors can be detected and corrected without destroying the information. Google demonstrated in 2023 that this type of code effectively reduces the error rate as the network grows.}, the Steane code, the Shor code~--- are devices that detect and repair coherence violations of a qubit without directly measuring its state (which would destroy it). In the language of the PDL, these are technological implementations of the principle of minimal leakage: they construct a second-level closure~--- a meta-closure~--- capable of absorbing the coherence leakage of the first-level closure (the physical qubit) without that leakage propagating into the computation. This reading has a structural implication: the fault-tolerance threshold\footnote{The fault-tolerance threshold is the error rate per operation below which quantum error correction is possible: if each operation has less than roughly 1\,\% error, one can in principle build an arbitrarily reliable quantum computer by adding correction qubits. Above this threshold, errors accumulate faster than they can be corrected.} of quantum computing should, within the PDL framework, be derivable from the leakage parameters $\varepsilon$ of the underlying architecture, rather than merely being calculated empirically.

\medskip

\textbf{Entanglement as a shared surface constraint.} Quantum entanglement is, in standard physics, a non-local correlation between two systems that cannot be described independently\footnote{Quantum entanglement was described by Einstein as ``spooky action at a distance'' (\textit{spukhafte Fernwirkung}): two entangled particles share a common state, such that measuring one instantaneously affects the description of the other, regardless of the distance separating them. This phenomenon has been confirmed experimentally many times, notably by Alain Aspect in 1982 and John Clauser, who were awarded the Nobel Prize in Physics in 2022.}. The PDL offers a structural interpretation: two entangled qubits are two closures that share a common active surface. Separating the qubits geographically does not break this constraint as long as the two closures continue to share a common boundary in the global relational network. This interpretation converges with the $\text{ER} = \text{EPR}$ conjecture\footnote{The $\text{ER} = \text{EPR}$ conjecture, proposed by Juan Maldacena and Leonard Susskind in 2013, suggests that two entangled particles (EPR, after Einstein, Podolsky, and Rosen) are connected by a geometric tunnel in space-time (ER, for Einstein--Rosen, i.e.\ a \textit{wormhole}). This conjecture would thus unify quantum mechanics and general relativity within a single framework.} of Maldacena and Susskind (2013), which proposes that the entanglement between two particles corresponds to the existence of a geometric bridge (\emph{wormhole}) between them. The PDL arrives at the same structure by an entirely different route: not through the geometry of space-times, but through the topology of signed graphs.

\medskip

\textbf{A testable prediction: the structure of quantum noise.} The PDL predicts that decoherence is not a purely random process: it is structured by the topology of the leaking closure. More precisely, the quantum noise spectrum~--- the frequency distribution of errors on a qubit~--- should reflect the topological structure of its environmental coupling, rather than being simply white or Markovian\footnote{A Markovian process is a memoryless process: the future state depends only on the present state, not on the past. Markovian noise is therefore temporally uncorrelated. Non-Markovian noise exhibits temporal correlations: errors at time $t$ are linked to past errors, reflecting an internal structure of the system.}. Recent experiments on superconducting quantum processors from IBM and Google (notably Sycamore and Heron) have revealed temporal correlations in the phase noise (\emph{$1/f$ noise})\footnote{$1/f$ noise (pink noise or flicker noise) is a type of noise whose power is inversely proportional to frequency: it is stronger at low frequencies than at high ones. This type of noise is ubiquitous in nature~--- found in transistor current fluctuations, heartbeats, and even music~--- and always signals the presence of long-memory processes.} that suggest precisely a non-Markovian structure. The PDL predicts that this structured noise should be modelisable as a hierarchical coherence leakage~--- with correlations at multiple temporal scales corresponding to the different levels of closure of the system. This model is computable and comparable to existing experimental data: it is one of the most immediately accessible tests of the programme.

\medskip

\textbf{Towards a PDL criterion for quantum advantage.} Quantum advantage~--- the capacity of a quantum computer to solve certain problems more efficiently than a classical computer~--- rests on sustaining coherence across a sufficient number of qubits for a sufficient length of time. In the PDL, this advantage has a precise structural interpretation: a quantum computation exploits the capacity of a multi-level closure to sustain superpositions of coherence configurations, whereas a classical computation can sustain only one configuration at a time. The PDL therefore predicts the existence of a \emph{closure threshold}~--- a minimal number of correlated logical qubits beyond which quantum advantage becomes irreversible~--- and that this threshold depends on the topology of the processor's coupling graph, not merely on its qubit count. This prediction is verifiable on current architectures: by comparing the performance of quantum processors with different topologies (linear chains, square grids, toroidal graphs\footnote{A toroidal graph is a network of connections whose topology is that of a torus~--- the surface of a doughnut. In such a network, every node has exactly the same number of neighbours and there are no ``edges'': the structure is periodic in all directions. Google uses a square-grid architecture; other architectures explore more complex topologies.}) at equal qubit numbers, the PDL predicts performance differences that are not explained by standard decoherence models.

\subsection*{Noetics: Consciousness, Cognition, and Coherence}

The PDL does not claim to explain consciousness. But it formulates a precise hypothesis about what consciousness \emph{would be} if the logic of coherence proved fundamental: a process of maintaining global coherence in a network of neuronal closures~--- that is, exactly what a complex physical system does in order to remain coherent in the face of perturbations. This hypothesis is refutable.

\medskip

\textbf{A prediction concerning neuronal decoherence.} If consciousness is a coherence phenomenon, then altered states of consciousness~--- sleep, anaesthesia, dissociative states~--- should correspond to measurable reductions in the functional coherence of the neuronal network. This is not in itself a new prediction: the Integrated Information Theory (IIT) of Tononi\footnote{Integrated Information Theory (IIT), proposed by neuroscientist Giulio Tononi, suggests that consciousness is proportional to the amount of information a system generates ``above and beyond'' what its parts generate separately~--- a measure called $\Phi$ (phi). The greater $\Phi$, the more conscious the system. This theory predicts that certain highly interconnected networks could be more conscious than a human brain, which remains highly controversial.} makes similar predictions. But the PDL predicts something more precise: the transition between consciousness and unconsciousness should have the structure of a \emph{closure phase transition}~--- not a continuous degradation of coherence, but an abrupt transition between a globally coherent closure regime and a regime of locally decoupled closures. High-resolution EEG\footnote{EEG (electroencephalography) measures the collective electrical activity of neurones through the skull. At high temporal and spatial resolution, it allows the functional coherence of the brain to be tracked in real time, particularly during transitions between states of consciousness.} measurements during anaesthetic induction~--- notably work on the ``loss of complexity'' measured by Lempel--Ziv entropy\footnote{Lempel--Ziv entropy is a measure of the compressibility of a signal: the more complex and unpredictable a signal, the less compressible it is and the higher its Lempel--Ziv entropy. Recent studies show that this measure drops sharply during anaesthetic induction, suggesting an abrupt loss of cerebral complexity~--- exactly what the PDL predicts for a closure phase transition.}~--- could test whether this transition is of first order (abrupt) or second order (gradual).

\medskip

\textbf{The paradox of Born's rule.} The PDL has sketched a derivation of Born's rule\footnote{Born's rule is the rule that allows probabilities to be calculated in quantum mechanics: if a system is in state $|\psi\rangle$, the probability of finding it in state $|\phi\rangle$ is $|\langle\phi|\psi\rangle|^2$. It is one of the most precisely verified rules in all of physics, but its deeper origin remains mysterious: why the squared modulus, and not the modulus itself, or its cube?}~--- the rule that predicts probabilities in quantum mechanics~--- from coherence constraints on active surfaces, without presupposing a priori the Hilbert-space structure of quantum mechanics\footnote{\textit{Relational Emergence of Born's Rule and a Golden-Ratio Active Surface in PDL}, v2, C.~Laubscher, 2026.}. If this derivation is correct, it implies that quantum probabilities are not fundamental properties of reality, but emergent properties of coherence management by observers that are themselves described as closures. This prediction has a testable consequence: in quantum interference experiments involving ``observers'' of very different sizes and complexities~--- atoms, molecules, silica beads, and ultimately biological systems~--- the PDL predicts that deviations from the standard Born rule could appear when the size of the observer approaches that of the observed system. These deviations would, according to the PDL, be a signature of coherence leakage between observer and system, and not a violation of quantum mechanics.

\subsection*{AI as an Accelerator of Testing}

An unexpected development of 2025--2026 deserves mention: the convergence between artificial intelligence and theoretical physics. Deep-learning models\footnote{Deep learning is an artificial intelligence technique that uses artificial neural networks with many layers to recognise patterns in massive datasets. These same architectures can be used to explore high-dimensional mathematical spaces, such as that of the PDL's admissible signed graphs.} are now capable of exploring high-dimensional combinatorial spaces~--- precisely the type of problem posed by the search for alternative architectures to $(24, 28, 930, 10\,087, 11\,017)$ within the space of admissible signed graphs. The second refutation scenario of the PDL~--- finding an alternative architecture reproducing $\alpha$ and $G$ with the same precision~--- is today computationally addressable: a search programme using gradient descent\footnote{Gradient descent is an optimisation algorithm: it seeks the minimum of a function by moving step by step in the direction of steepest descent. Applied to the space of the PDL's integer architectures, it would allow one to systematically explore whether other combinations of integers satisfy the same coherence constraints.} in the space of integers constrained by the PDL's coherence criteria is technically feasible on current computing infrastructure. If such an exhaustive search finds no alternative within a reasonably large space, the PDL's uniqueness conjecture emerges considerably strengthened. If it finds one, the programme must reformulate itself. In either case, the answer is accessible: this is a programme of computational research, not merely theoretical.

\subsection*{What Remains to Be Built}

An honest scientific programme does not merely enumerate its predictions: it identifies the open construction sites, those without which predictions will remain conjectures.

\begin{itemize}
\item \textbf{A complete derivation of the electron-electron scattering cross-section from the $(4,6)$ structure.} Without this calculation, the PDL cannot be confronted with LHC data at high energies.
\item \textbf{A calculation of the cosmological constant $\Lambda$ from the number of elementary closures in the observable universe.} This is the most direct cosmological test; it is also the most difficult.
\item \textbf{A complete formalism for discrete quantum gravity.} The PDL sketches an emergent metric via coherence cycles and an analogue of Dirac spinors via the $(4,6)$ structures\footnote{\textit{Towards an Einstein--Dirac Unification from the PDL Framework}, C.~Laubscher, 2026.}. The strong-field limit~--- the regime of black holes and early-universe cosmology~--- has not yet been treated.
\item \textbf{A PDL criterion for biological viability.} Synthetic biology awaits a structural criterion: what is the necessary and sufficient condition, in terms of closures, for a biological system to be viable? This criterion remains to be formulated.
\item \textbf{A formal correspondence between PDL structures and quantum computing protocols.} The preceding section lays the conceptual foundations; the translation into the language of quantum circuits and operator algebra remains to be done.
\end{itemize}

These construction sites are not lacunae: they are the proof that the programme is alive. A programme that has nothing left to do is no longer a research programme; it is a museum. The PDL, at this stage, is far more a work in progress than a completed edifice~--- and that is precisely what makes it a serious programme.

% ============================================================
\backmatter
\appendix

\chapter*{Appendix — Going Further}
\addcontentsline{toc}{chapter}{Appendix}

\subsection*{A. Glossary of Key Terms}

\begin{description}

\item[Active surface] Subset of the relations of a composite structure that is projected outward and available for couplings with other structures.

\item[Coherence budget] The total amount of internal coherence that a structure can maintain simultaneously. Mass, energy, and certain physical constants are ways of measuring this budget at different scales.

\item[Coherence leakage ($\varepsilon$)] The irreducible fraction of triangles that cannot be made coherent simultaneously in a composite structure. This leakage is exported to the environment and constitutes, on large scales, the basis of gravitational coupling.

\item[Coherence ratio] Way of reading fundamental physical constants in the PDL: not as arbitrary numbers, but as measures of how coherence is distributed between the interior, the surface, and the exterior of a given structure.

\item[Closure] A relational structure that satisfies its own coherence constraints without appealing to an external support. To exist, in the PDL, is to instantiate a closure.

\item[Fine-structure constant ($\alpha$)] Measures the strength of electromagnetic coupling. In the PDL, it is reinterpreted as the ratio between the active surface of the proton and its total relational budget.

\item[Finite signed graph] A finite set of nodes linked by edges, each carrying a sign $+1$ or $-1$. This is the basic substrate on which all PDL closures are defined.

\item[Hierarchical filtering] Mechanism by which tiny coherence leakage at the nuclear scale is amplified, level by level, until it produces an observable macroscopic effect~--- notably gravitation.

\item[Minimal stationary closure] The smallest structure that can both pulse and autonomously maintain its triangular coherence. In the PDL, this is the $(4,6)$ structure, identified with the logical archetype of the electron at rest.

\item[Projective Dynamic Logo (PDL)] Fundamental research programme developed by Cédric Laubscher since 2024, which reconstructs physical reality from axioms defined on finite signed graphs, without presupposing space, time, or particles.

\item[Pulse (binary pulsation)] Alternation between two distinct states that returns to itself in two steps. This is the minimal form of existence in the PDL: a difference capable of repeating itself.

\item[Quantum entanglement] A property of two quantum particles that have interacted, whose states can no longer be described independently of one another, even when separated by a large distance. Measuring one instantaneously affects what can be predicted for the other, without any physical signal being exchanged between them. This phenomenon, predicted by quantum mechanics and confirmed experimentally since the 1980s (Alain Aspect's experiments), constitutes one of the deepest differences between quantum and classical physics. In the PDL, entanglement is reinterpreted as a coherence shared between the boundaries of two closures, via mixed triangles: it is not a mysterious action at a distance, but the relational signature of a structural constraint that two closures hold in common.

\item[Relational sea] In the proton architecture, the set of dynamic links that surround the stationary cores. It is analogous to the gluon and virtual-quark field of quantum chromodynamics.

\item[(4,6) structure] Complete graph on four vertices and six edges: the first logically admissible closure under the PDL constraints. Identified with the archetype of the electron at rest.

\item[Transcendence (in the PDL sense)] The structural property of any sufficiently complex closure: the impossibility of fully containing itself, the necessity of depending on a global coherence field that exceeds its own boundaries.

\item[Triangular coherence] Condition imposed on each triplet of mutually connected relations: their product must satisfy a precise sign rule. A coherent triangle is a closed loop without internal contradiction.

\end{description}

\subsection*{B. Map of Technical Publications}

All technical work on the PDL is freely available on \href{https://zenodo.org}{\textbf{Zenodo}}, under the name Cédric Laubscher. The following publications constitute the core of the programme:

\begin{enumerate}

\item \textit{The Emergence of Physical Reality within the Framework of the Projective Dynamic Logo (PDL)} — axioms, derivation of the (4,6) structure, correspondence principle with the electron.

\item \textit{Minimal Stationary Closures in the PDL Framework: Necessity of the (4,6) Block} — rigorous demonstration of the minimality of the (4,6) structure.

\item \textit{On the Combinatorial Selection of the Proton Architecture in PDL} — derivation of the integers 24, 28, 930, 10\,087, 11\,017 and uniqueness conjecture.

\item \textit{A Structural Derivation of the Fine-Structure Constant from the PDL} — combinatorial expression for $\alpha$ without an adjustable parameter.

\item \textit{Coherence Leakage, Hierarchical Filtering, and the Exponent 18 in the PDL Framework} — leakage mechanism and amplification to effective gravitational coupling.

\item \textit{From Discrete Coherence Flux to Effective Fields} — passage from discrete graphs to continuous field descriptions.

\item \textit{Discrete Cavity Modes and Logical Density of States in the PDL Framework} — standard density of states reproduced from finite graphs.

\item \textit{Logical Leakage as Self-Maintained Probability: A Topological Reformulation in the PDL Framework} — quantum probabilities reinterpreted as frequencies of coherence violations.

\item \textit{Towards a Schrödinger-Type Dynamics from the PDL} — link between coherence cycles and the Schrödinger equation.

\item \textit{Towards an Einstein--Dirac Interface in the PDL Framework} — bridges between the PDL, general relativity, and relativistic quantum mechanics.

\item \textit{PDL: Global Mapping of Structures, Results, and Open Problems} — overview of the programme, established results, conjectures, and open questions.

\item \textit{Whoever We May Be: Existence, Coherence, and Vulnerable Instruments} — ontological and philosophical synthesis of the PDL corpus for a mixed audience.

\end{enumerate}

% ============================================================
\end{document}